\section{Retained rank data and an exact carrier}\label{sec:high-preparation}

Our aim is to retain the source energy carried by the regular rank-one
and rank-two points while preparing most of the target for exact local
replacements.  The retained data will first be flattened onto separated
plates, with only a small change to their derivatives.  On opposite
sides of each plate we select paired target rectangles.  The carrier
of Proposition~\ref{hp:carrier} will agree with affine or PL
homeomorphisms over neighborhoods of these rectangles and over an open
set with arbitrarily small residual target volume, while agreeing with
the chosen map near the retained source data.  The order of
these two requirements explains its two rounds: endpoint neighborhoods
are fixed first; the remaining target volume is prescribed afterwards.

The carrier may collapse isolated tame balls.  Such fibres permit the
energy-controlled local replacements below and can later be opened
topologically.  For $2<p<3$, the replacements can also leave marked
source interiors whose energy is treated by a separate compression.
The final two subsections prepare that compression and the eventual
full-rank correction used in Section~\ref{sec:high-assembly}.

\subsection{Regular points and fine tolerances}

Throughout this section, $p>2$ and the topological orientation of $f$
is positive.  A reflection reduces the other
orientation to this one.  All regularity assertions below are local on
the open sets; no boundary regularity is involved.  A map has the volume
Lusin property $N$ if it maps sets of three-dimensional Lebesgue
measure zero to sets of measure zero.

We use the weighted area formula in the form recorded by
Simon~\cite[Chapter~2, Section~3, Equations~(3.4)--(3.5)]{Simon2014Draft}.
The general Lipschitz-decomposition background for Sobolev maps is
provided by Bojarski--Haj\l asz~\cite[Theorem~3(1)]{BojarskiHajlasz1993}.
The stronger ordinary differentiability and image-null assertions
needed here are proved next, using openness and two-dimensional
Sobolev estimates on spheres.

\begin{lemma}[Regular points of a Sobolev homeomorphism]\label{hp:regularity}
Let $u:G\to G'$ be a homeomorphism between open subsets of $\R^3$,
with $u\in W^{1,p}_{\mathrm{loc}}(G;\R^3)$ and $p>2$.
There is a Borel set $R_u\subset G$ of full measure such that, for
$x\in R_u$, the weak differential $A=Du(x)$ is also the Fr\'echet
differential and
\begin{equation}\label{hp:regular-point-tests}
 \lim_{r\downarrow0}\sup_{0<|z-x|\le r}
 \frac{|u(z)-u(x)-A(z-x)|}{|z-x|}=0,
 \qquad
 \lim_{r\downarrow0}\frac{1}{|B_r|}
 \int_{B_r(x)}|Du-A|^p=0.
\end{equation}
If $u$ has positive topological orientation, then $\det A>0$ at
all points of $R_u$ where $A$ is invertible.  Moreover,
\begin{equation}\label{hp:deficient-image-null}
 \cL^3\bigl(u(\{x\in R_u:\rank Du(x)<3\})\bigr)=0.
\end{equation}
If $p\ge3$, $u$ has the volume Lusin property $N$ locally, and hence
on $G$.  Finally, for every fixed coordinate direction, almost every
parallel coordinate-plane section of $G$ has image under $u$ of
three-dimensional measure zero.  This last assertion holds for every
continuous $W^{1,p}_{\mathrm{loc}}$ map, without injectivity.
\end{lemma}

\begin{proof}
We first pass from Sobolev differentiation in mean to ordinary
Fr\'echet differentiation.  At almost every $x$, with $A=Du(x)$,
\begin{equation}\label{hp:rescaled-jet-small}
 w_r(\xi):=\frac{u(x+r\xi)-u(x)}{r}-A\xi
 \longrightarrow0 \quad\hbox{in }W^{1,p}(B_2;\R^3).
\end{equation}
For completeness, the gradient assertion is the Lebesgue differentiation
Theorem applied to $Du$.  To obtain the value assertion, subtract the
affine function with gradient $A$ and apply Poincar\'e's inequality on
concentric balls.  The difference between its mean on a ball of radius
$t$ and on the concentric ball of radius $t/2$ is bounded by
\[
 Ct\left(\frac{1}{|B_t|}\int_{B_t(x)}|Du-A|^p\right)^{1/p}.
\]
Summing on dyadic radii and using the precise value $u(x)$ gives the
value assertion in \eqref{hp:rescaled-jet-small}.

Fix $0<d<1/8$.  For every $z\in B_1$, slicing the annulus
$B_{2d}(z)\setminus B_d(z)$ gives a radius $\rho\in(d,2d)$ for which
\[
 \int_{\partial B_\rho(z)}
   (|w_r|^p+|D_\tau w_r|^p)
 \le d^{-1}\|w_r\|_{W^{1,p}(B_2)}^p.
\]
The trace agrees with the continuous restriction.  Sobolev embedding
on the two-dimensional sphere, using $p>2$, therefore gives
\[
 \sup_{\partial B_\rho(z)}|w_r|
 \le C_{p,d}\|w_r\|_{W^{1,p}(B_2)}.
\]
The constant is independent of $z$.  Each coordinate of the
homeomorphism
$u_r(\xi)=(u(x+r\xi)-u(x))/r$ attains its maximum and minimum on
the boundary of $\overline B_\rho(z)$: an extremum in the interior
would contradict openness of $u_r$.  Comparing the affine function
$A\xi$ at $z$ and on this sphere yields
\[
 \sup_{z\in B_1}|w_r(z)|
 \le C\|A\|_{\op}d+C_{p,d}\|w_r\|_{W^{1,p}(B_2)}.
\]
First let $r\downarrow0$, then $d\downarrow0$.  We obtain uniform
affine approximation on $B_r(x)$, which is equivalent to the first
assertion of \eqref{hp:regular-point-tests}.  Intersecting the full
measure sets used here gives a Borel choice of $R_u$.

If $A$ is invertible, uniform approximation on a sufficiently small
sphere makes $u(x+r\xi)-u(x)$ homotopic, through maps avoiding zero,
to $rA\xi$.  The local degree of the homeomorphism is therefore
$\operatorname{sgn}\det A$, proving the orientation assertion.

The set of Fr\'echet differentiability points admits a countable
Lipschitz decomposition.  Indeed, restrict first to points where
$|Du|\le m$ and
\[
 |u(z)-u(x)-Du(x)(z-x)|\le |z-x|
 \quad(0<|z-x|<1/m),
\]
and then partition into sets of diameter less than $1/m$ in an
interior exhaustion.  On each resulting set the restriction of $u$
is Lipschitz.  Such sets cover every differentiability point after
varying $m$.  On a Lipschitz extension of each restriction, its
approximate differential agrees almost everywhere on that set with
$Du$.  The area formula, applied to the subset where $\rank Du<3$,
now proves \eqref{hp:deficient-image-null}.  Source null subsets of
these Lipschitz pieces have null images as well, so the conclusion
concerns the entire regular deficient-rank set, not merely a source
full-measure subset of it.

For the volume property $N$, it suffices to use $p=3$ locally.
For an interior cube $Q$, choose a sphere of radius comparable to
its side length, enclosing $Q$ and contained in a fixed enlargement
$C_0Q\Subset G$.  The sphere oscillation inequality and slicing give
\[
 (\operatorname{osc}_Q u)^3
 \le C\int_{C_0Q}|Du|^3.                 \tag{$*$}
\]
Here openness again bounds coordinate oscillations inside the sphere
by those on its boundary.  If $E\Subset G$ is null, choose an open
neighborhood $O\Subset G$ of $E$ with arbitrarily small
$\int_O|Du|^3$.  A Whitney decomposition of $O$ may be chosen with
$C_0Q\subset O$ and bounded overlap of these enlargements.  Summing
$(*)$ over the cubes that cover $E$ bounds the outer volume of $u(E)$
by $C\int_O|Du|^3$.  This proves the claim by exhaustion.  When
$p>3$, local $L^3$ integrability follows from H\"older's inequality.

For the last assertion, Fubini gives a continuous $W^{1,p}$
restriction to almost every coordinate-plane section.  On a compact
square in such a section, divide into squares $Q$ of side length
$a$, allowing uniformly bounded enlargements inside a slightly larger
compact square.  Planar Morrey's inequality and H\"older's inequality
give
\[
 \begin{split}
 \sum_Q\diam(u(Q))^2
 &\le C\sum_Q
 a^{2-4/p}\left(\int_{2Q}|D_\tau u|^p\right)^{2/p}\\
 &\le C\left(\sum_Q a^2\right)^{1-2/p}
       \left(\sum_Q\int_{2Q}|D_\tau u|^p\right)^{2/p}.
 \end{split}
\]
This bound is independent of $a$.  Uniform continuity gives
\[
 \max_Q\diam(u(Q))\longrightarrow0,
 \qquad \sum_Q\diam(u(Q))^3\longrightarrow0.
\]
A countable exhaustion proves the three-dimensional null-image
assertion.
\end{proof}

\begin{remark}[Fine value tolerances]\label{hp:fine-tolerance}
Lemma~\ref{pre:local-tolerances} supplies, for every positive continuous
function $a$ on an open set $G$ and every $L>0$, a positive
$L$-Lipschitz minorant $a_0\le a$ with
$a_0(x)\le L\dist(x,\R^3\setminus G)$.  We use these minorants to
combine the fine value requirements below, including distance to the
target boundary and previously fixed compact-set accuracies.
\end{remark}


\subsection{Flattening the deficient positive ranks in the target}

The regular deficient-rank image has zero target volume, but the
integral of $|Df|^p$ over its source preimage need not be small.  We
therefore organize this part of the map rather than discard it.
The next construction moves almost all of its weighted data onto
separated flat plates and preserves their tangential derivatives.

Let $R_f$ be the regular set from Lemma~\ref{hp:regularity}, and set
\[
 D=\{x\in R_f:1\le\rank Df(x)\le2\},
 \qquad
 \mu(E)=\int_{E\cap D}(1+|Df|^p)\,dx,
 \qquad
 \nu=f_\#(\mathbf1_\Omega\,\cL^3).
\]
The finite measure $\mu$ is the measure used for all losses in this
preparation.  Boundedness of $\Omega$ and $f\in W^{1,p}$ ensure
its finiteness.

The decomposability bundle $V(\nu,y)$ records the tangent directions
of rectifiable curves whose averaged measures are absolutely continuous
with respect to $\nu$.  We use the precise locality and differentiability
properties cited in the proof below to find a direction transverse to
this bundle on the retained data.

\begin{lemma}[A forbidden cone for the transported measure]
\label{hp:transported-bundle}
Let $V(\nu,y)$ be the decomposability bundle of $\nu$.  Then
\begin{equation}\label{hp:range-in-bundle}
 \operatorname{im}Df(x)\subset V(\nu,f(x))
 \quad\hbox{for almost every }x\in\Omega.
\end{equation}
Moreover, $V(\nu,y)\ne\R^3$ for $\nu$-almost every $y\in f(D)$.
Consequently, for each $0<\lambda<1$, $f(D)$ can be partitioned,
up to a $\nu$-null set, into finitely many measurable pieces with
associated unit vectors $e$ such that
\begin{equation}\label{hp:almost-normal-direction}
 |\operatorname{proj}_{V(\nu,y)}e|<\lambda/10
 \quad\hbox{on the corresponding piece}.
\end{equation}
In coordinates $(h,v)\in e^\perp\times\R$, the bundle on that piece
contains no nonzero vector in the double cone
\[
 \{(a,b):|a|\le |b|/\lambda\}.
\]
\end{lemma}

\begin{proof}
We use the decomposability bundle and its strong locality property
from Alberti--Marchese \cite[Section~2.6 and Proposition~2.9(i)]
{AlbertiMarchese2016}.  In particular, tangent directions of measurable
families of rectifiable curve measures belong to $V(\nu,\cdot)$
whenever the aggregate curve measure is absolutely continuous with
respect to $\nu$.  Weighted restrictions are allowed: a density can
be absorbed by restriction and the layer-cake formula.

Fix a coordinate direction $e_i$ and an interior rectangular box.
For almost every line in that direction, the restriction of $f$ is
absolutely continuous.  It is also injective, so the one-dimensional
area formula identifies the pushforward of
$|\partial_i f|\,dt$ with length measure on the image curve.
Integration over the transverse coordinates gives the aggregate
measure
\[
 f_\#(\mathbf1_{\text{box}}|\partial_i f|\,dx)\ll\nu.
\]
The images are rectifiable for almost every line; if necessary they
may be split into bounded-length pieces or parametrized by arclength.
At almost every parameter where $\partial_i f\ne0$, their tangent
line is $\operatorname{span}\{\partial_i f\}$.  The defining property
of the bundle therefore gives
$\partial_i f(x)\in V(\nu,f(x))$ for almost every $x$ in the box
where that derivative is nonzero.  A countable box exhaustion and the
three coordinate directions prove \eqref{hp:range-in-bundle}.

By Lemma~\ref{hp:regularity}, $f(D)$ has zero volume.  Suppose that
$A\subset f(D)$ has positive $\nu$ measure and
$V(\nu,y)=\R^3$ on $A$.  The differentiability Theorem of
Alberti--Marchese \cite[Theorem~1.1(i)]{AlbertiMarchese2016} says that
every Lipschitz function on $\R^3$ is ordinarily differentiable
$\nu\vert_A$-almost everywhere.  The converse to Rademacher's Theorem
of De Philippis--Rindler \cite[Theorem~1.14]{DePhilippisRindler2016}
then implies $\nu\vert_A\ll\cL^3$, which is impossible since
$A\subset f(D)$ is null.  This proves properness.  This invocation
uses the measure-theoretic converse to Rademacher, rather than an
uncited inference from the existence of individual curve families.

Choose a finite sufficiently fine net of directions on the unit
sphere.  Every proper linear subspace has a unit normal, so one
of these directions satisfies \eqref{hp:almost-normal-direction}.
Assign the first such direction measurably.  If $(a,b)$ belongs to
the displayed cone and is nonzero, then
\[
 \frac{|b|}{\sqrt{|a|^2+b^2}}
 \ge\frac{\lambda}{\sqrt{1+\lambda^2}}>\lambda/2,
\]
whereas any vector $w\in V(\nu,y)$ obeys
$|e\cdot w|< (\lambda/10)|w|$.  The cone is therefore forbidden.
\end{proof}

The simultaneous null-set conclusion below is a graph-parametrized form of
Alberti--Marchese's Lemma~7.5.  Its compact-convex selection mechanism
is Rainwater's lemma, as presented in
\cite[Lemma~7.1 and Corollary~7.2]{AlbertiMarchese2016}.
We give the argument because one common exceptional set for every
allowed graph is essential to the subsequent flattening.

\begin{lemma}[A simultaneous graph-null set]\label{hp:cone-null}
Fix orthogonal coordinates $(h,v)\in\R^2\times\R$ and
$0<\lambda<1$.  Suppose a finite measure $\sigma$ is concentrated on a
Borel set $A$ and is absolutely continuous with respect to $\nu$,
and that $V(\nu,y)$ contains no nonzero vector in
$\{(a,b):|a|\le |b|/\lambda\}$ for $\sigma$-almost every $y$.
There is a Borel set $A_0\subset A$ of full $\sigma$ measure such that
\begin{equation}\label{hp:graph-null}
 \bigl|\{v:(h(v),v)\in A_0\}\bigr|=0
\end{equation}
for every $1/\lambda$-Lipschitz graph $h$ on any interval.
\end{lemma}

\begin{proof}
We first work in a compact rectangle $X=H\times I$, where $H$ is a
closed horizontal box and $I$ a compact vertical interval.  Let
$\Gamma$ be the compact family of all $1/\lambda$-Lipschitz maps
$h:I\to H$, equipped with the uniform topology.  Write
\[
 \rho_h=(v\mapsto(h(v),v))_\#(\cL^1\vert_I),
 \qquad
 \mathcal C=\left\{\int_\Gamma\rho_h\,dP(h):
                         P\in\mathcal P(\Gamma)\right\}.
\]
The map $h\mapsto\rho_h$ is weakly continuous.  Thus $\mathcal C$
is convex and weakly compact in the finite measures on $X$.

Every $\rho\in\mathcal C$ is singular with respect to
$\sigma\vert_X$.  Indeed, otherwise there would be a nonzero measure
$\sigma_0\le\rho$ with $\sigma_0\ll\sigma\vert_X$.  Writing
$a=d\sigma_0/d\rho$ and $\rho=\int\rho_h\,dP(h)$ yields
\[
 \sigma_0=\int_\Gamma a\rho_h\,dP(h).
\]
Each $\rho_h$ is equivalent to length measure on its graph, with
bounded positive density.  Density restriction and layer-cake
therefore express $\sigma_0$ as a family of rectifiable curve measures.
Its aggregate is absolutely continuous with respect to $\nu$.
The tangents of the graphs must belong to $V(\nu,\cdot)$ at almost
every point charged by this family, while their directions
$(h'(v),1)$ lie in the forbidden cone.  This contradicts
$\sigma_0\ne0$.

For $0\le\phi\le1$ continuous on $X$, put
\[
 L(\phi,\rho)=\int_X\phi\,d\rho
                 +\int_X(1-\phi)\,d\sigma.
\]
Mutual singularity and regularity of finite measures imply
$\inf_\phi L(\phi,\rho)=0$ for every $\rho\in\mathcal C$.
The compact-convex minimax Theorem, with $\mathcal C$ as its compact
side \cite[Theorem~4.2]{Sion1958}, applies to this continuous affine functional
and gives
\[
 \inf_{\substack{\phi\in C(X)\\0\le\phi\le1}}
     \sup_{\rho\in\mathcal C}L(\phi,\rho)=0.
\]
Choose $\phi_n$ with
\[
 \int_X(1-\phi_n)\,d\sigma
       +\sup_{h\in\Gamma}\int_I\phi_n(h(v),v)\,dv
 \le2^{-n}.
\]
Tonelli's Theorem shows that $\phi_n\to1$ at $\sigma$-almost every
point, while, for each individual graph, $\phi_n(h(v),v)\to0$
for almost every $v$.  Hence the single Borel set
$\{y:\phi_n(y)\to1\}$ has full $\sigma\vert_X$ measure and is
null on every graph in $\Gamma$.  No intersection of graph-dependent
full-measure sets is being taken.

To pass to arbitrary graphs, use countably many rectangles with
strict interior margins and exhaust the measure by their interiors.
On the vertical interval of such a rectangle, restrict the graph
and project its horizontal component to the closed box $H$; this
preserves the Lipschitz constant and does not change graph points
inside the smaller rectangle.  A graph given on a shorter interval
extends with the same Lipschitz constant by constant extension at its
endpoints.  Taking a countable union of the resulting full-measure,
locally graph-null sets proves \eqref{hp:graph-null}.  Intersect this
set with $A$ at the end.
\end{proof}

We next convert graph-nullness into finitely many thin bands.  The
cone order, antichain decomposition and Lipschitz slabs have their
counterparts in the null-set constructions of
Alberti--Cs\"ornyei--Preiss~\cite[Sections~2.2 and~3]{AlbertiCsornyeiPreiss2010}.
The uniform near-hit estimate is the compactness mechanism used in
\cite[Lemma~4.12, Step~1]{AlbertiMarchese2016}; we retain its proof here.

\begin{lemma}[Finite thin bands]\label{hp:thin-bands}
Let $K$ be a compact subset of the interior of a rectangular chart
$H\times I\subset\R^2\times\R$, and assume that $K$ satisfies
\eqref{hp:graph-null}.  For every $\delta>0$, there are finitely many
piecewise affine functions $g_1,\dots,g_N$ on $H$, each with
Lipschitz constant at most $C\lambda$, and a common length $L>0$
such that
\[
 NL<\delta,
 \qquad
 K\Subset\bigcup_{j=1}^N
       \{(h,v):|v-g_j(h)|\le L/2\}.
\]
The values of the functions, and their bands, can be kept in a fixed
compact subinterval of the interior of $I$.  Here the containment
$\Subset$ means that the union contains an open neighborhood of $K$.
\end{lemma}

\begin{proof}
Take a cubical grid of spacing $w$ and let $P_w$ be the centers of
the grid boxes that meet $K$.  For distinct centers define
\[
 (h,v)\prec(h',v')\quad\Longleftrightarrow\quad
                         v'-v\ge\lambda|h'-h|.
\]
This is a strict partial order: transitivity follows from the
triangle inequality, and distinct comparable centers have $v'>v$.
Because vertical grid coordinates differ by integer multiples of
$w$, successive centers in a chain have vertical separation at least
$w$.

The horizontal coordinates of a chain, interpolated linearly as a
function of $v$ and extended constantly beyond its first and last
nodes, form an allowed $1/\lambda$-Lipschitz graph.  All such graphs
lie in a fixed compact graph family, after slightly enlarging $H$
and $I$ while keeping the chart margins.  For this family set
\[
 \omega(r)=\sup_{h\in\Gamma}
  \bigl|\{v\in I:\dist((h(v),v),K)\le r\}\bigr|.
\]
Then $\omega(r)\to0$ as $r\downarrow0$.  Otherwise compactness
would give uniformly convergent graphs along a sequence $r\downarrow0$;
the upper limit of the near-hit indicators is bounded by the
indicator that the limiting graph meets $K$.  Reverse Fatou and
\eqref{hp:graph-null} give a contradiction.

Around each node of a chain take a vertical interval of radius
$w/4$.  These intervals are disjoint and, for small $w$, lie inside
the fixed interval with margins.  On each such interval the graph
is within $C_\lambda w$ of $K$.  If $n_w$ is the longest chain
length, it follows that
\[
 \tfrac12 n_ww\le\omega(C_\lambda w),
 \qquad\hbox{hence}\qquad n_ww\longrightarrow0.
\]
Give every center its longest-chain rank.  Centers of the same rank
form an antichain, and two of them satisfy
$|v'-v|<\lambda|h'-h|$.  In particular, their horizontal coordinates
are distinct and their heights extend to a $\lambda$-Lipschitz
function on $H$ by the scalar Lipschitz extension formula.  Clip
the extension to an interior vertical interval containing all the
center heights.  Piecewise affine interpolation on a sufficiently
fine regular triangulation then approximates it to $o(w)$ with
Lipschitz bound $C\lambda$.

There are at most $n_w$ such functions.  A band of length $L=C_2w$
about each function contains the boxes assigned to it with a strict
margin, provided $C_2$ is fixed sufficiently large.  Thus the bands
contain a neighborhood of $K$, and
$NL\le C_2n_ww\to0$.  The initial margins and a sufficiently small
$w$ give the final assertion.
\end{proof}

\begin{proposition}[Target flattening]\label{hp:flattening}
Given $\eps>0$, $0<\lambda<1/100$, and a positive continuous fine
tolerance $a$ on $\Lambda$, one can retain a compact set $K\subset D$
with $\mu(D\setminus K)<\eps$ and construct a Lipschitz map
$S_0:\Lambda\to\Lambda$ with the following properties.
\begin{enumerate}
\item The map is the identity outside finitely many disjoint
rectangular boxes compactly contained in $\Lambda$,
$\Lip(S_0)\le C$, and $|S_0(y)-y|<a(y)$.  The constant $C$ is
absolute.  For $0<\kappa\le1$,
\[
 S_\kappa=(1-\kappa)S_0+\kappa\Id
\]
is a bi-Lipschitz homeomorphism of $\Lambda$ onto itself.
\item With $F_\kappa=S_\kappa\circ f$ for $0\le\kappa\le1$,
\begin{equation}\label{hp:flattening-derivative}
 |DF_\kappa|\le C|Df|\quad\hbox{a.e. on }\Omega,
 \qquad
 |DF_\kappa-Df|\le C\lambda|Df|
                   \quad\hbox{a.e. on }K.
\end{equation}
Moreover, $F_\kappa\to F_0$ strongly in $W^{1,p}(\Omega;\R^3)$,
and $\rank DF_0=\rank Df$ almost everywhere on $K$.
\item The compact set $F_0(K)$ is contained in finitely many compact
flat plate patches with pairwise disjoint ambient neighborhoods.
Each patch lies in a plane of slope at most $C\lambda$ in one of
the chosen coordinates.  Each patch has an assigned rank, either
one or two; the rank-one patches can in addition retain any prescribed
small directional grouping of $\operatorname{im}Df$.
\item There is a compact zero-volume set $Z\Subset\Lambda$,
contained in finitely many piecewise affine graph sheets, such that
\begin{equation}\label{hp:flattening-homeomorphic-complement}
 S_0:\Lambda\setminus S_0^{-1}(Z)\longrightarrow\Lambda\setminus Z
\end{equation}
is a locally bi-Lipschitz homeomorphism.  Its only non-singleton
fibres are the collapsed vertical intervals over points of $Z$.
Consequently $F_0$ is a homeomorphism from
$\Omega\setminus F_0^{-1}(Z)$ onto $\Lambda\setminus Z$.
\end{enumerate}
For almost every target point outside $Z$, its unique preimage under
this last homeomorphism either is a regular point of $F_0$ with
positive invertible differential and the power-mean test in
\eqref{hp:regular-point-tests}, or belongs to a fixed source null set.
When $p\ge3$, the second alternative occurs only on a target null
set.
\end{proposition}

\begin{proof}
We describe the measure selections before the geometric construction.
Since $\mu$ is finite, first discard an arbitrarily small $\mu$-mass so
that the remaining jets and the reciprocals of their nonzero singular
values are bounded.  Partition by rank and, at rank one, by a finite
angular partition of the range lines.  Lemma~\ref{hp:transported-bundle}
further partitions the data among finitely many directions $e$.
The measure $f_\#\mu$ is absolutely continuous with respect to $\nu$,
so every $\nu$-full choice in these pieces is also full for the
weight being retained.  Apply Lemma~\ref{hp:cone-null} to obtain
simultaneous graph-null data in each direction.

By inner regularity choose compact subsets of the finitely many
pieces with small total loss.  Their images are disjoint compact
sets, hence have positive separation.  In each assigned coordinate
frame take a sufficiently fine box grid inside that separation
margin and inside $\Lambda$.  Its cut planes can be chosen to have
zero $f_\#\mu$ measure: for each coordinate there are only countably
many levels with positive mass.  Discard small neighborhoods of the
finitely many cuts and take compact subsets again.  We have reduced
the problem, with arbitrarily small loss, to finitely many disjoint
boxes, with the compact data strictly inside each box and a fixed
rank and direction assignment there.  In what follows $K^*$ denotes
the compact target data in one such box.

Apply Lemma~\ref{hp:thin-bands} to $K^*$.  Sort the starts of the
length-$L$ bands pointwise, writing them as
$a_1(h)\le\cdots\le a_N(h)$.  Order statistics of finitely many
piecewise affine functions are piecewise affine after finite
subdivision, and retain their common Lipschitz bound.  Define
\begin{equation}\label{hp:separated-band-starts}
 s_1=a_1,\qquad s_j=\max\{a_j,s_{j-1}+L\}\quad(2\le j\le N).
\end{equation}
These starts remain $C\lambda$-Lipschitz, since equivalently
$s_j=\max_{i\le j}(a_i+(j-i)L)$.  Their bands have disjoint
interiors in each vertical column.

No point of an original band is lost.  To verify this, group the
successive new intervals into maximal contiguous clusters.  A cluster
starts at an unchanged start $s_i=a_i$, and every advanced interval
is attached to the top of its preceding cluster.  Since the original
starts are sorted, the portion below an advanced start belongs to
that preceding cluster; the remaining portion is covered by the
advanced interval itself.  This proves coverage by induction.
The total upward displacement is at most $NL$, so choosing $NL$
small preserves all vertical margins.

Choose a horizontal cutoff $0\le\chi\le1$, compactly supported in
the horizontal interior of the box and equal to one near the
horizontal projection of $K^*$.  Choose a nondecreasing Lipschitz
function $\eta(v)$ that is zero below a fixed level above every band
and is one near the top of the box.  Its transition is contained
strictly above all bands.  The term involving $\eta$ compensates
above the bands for the vertical length removed below it, so each
column keeps its endpoints.  For
$[t]_{[0,L]}=\max\{0,\min\{t,L\}\}$, define
\begin{equation}\label{hp:triangular-collapse}
 q_0(h,v)=v+\chi(h)\left(
       -\sum_{j=1}^N[v-s_j(h)]_{[0,L]}+NL\eta(v)\right),
 \qquad S_0(h,v)=(h,q_0(h,v)).
\end{equation}
Choose $NL$ after the cutoffs so small that
\begin{equation}\label{hp:band-cutoff-smallness}
 NL\Lip(\chi)\le\lambda,
 \qquad NL\Lip(\eta)\le1,
\end{equation}
and so that $NL$ is below the prescribed fine motion tolerance and
all unused box margins.

The correction in \eqref{hp:triangular-collapse} vanishes below all
bands, above the transition of $\eta$, and near the horizontal
boundary.  Thus $S_0$ equals the identity near the boundary of the
box.  For fixed $h$, the vertical map is nondecreasing, with slope
$1-\chi(h)$ on each band and slope
$1+\chi(h)NL\eta'(v)$ on the upper compensation interval.  Elsewhere
its slope is one.  These statements hold almost everywhere and
imply the corresponding monotonicity for the Lipschitz function.
In particular the slope lies in $[0,2]$.

The horizontal bound is independent of the number of bands.  At any
point of a piecewise affine region at most one truncated summand in
\eqref{hp:triangular-collapse} is unsaturated.  Every fully saturated
summand is the constant $L$.  The horizontal gradient of the sum
therefore has norm at most $C\lambda$.  Since the whole correction
has size at most $NL$, the cutoff term is bounded by
$NL\Lip(\chi)$.  Continuity across the finitely many pieces yields
\begin{equation}\label{hp:anisotropic-correction}
 |R(h,v)-R(\widetilde h,\widetilde v)|
 \le C\lambda|h-\widetilde h|+C|v-\widetilde v|,
 \qquad R(h,v)=q_0(h,v)-v.
\end{equation}
Thus $\Lip(S_0)\le C$ and $|S_0-\Id|\le NL$.  Each vertical
column is mapped onto itself, so the whole box is mapped onto itself.
Use this formula in each of the finitely many disjoint boxes and the
identity elsewhere.  Extension by the identity to $\R^3$ is globally
Lipschitz, as is seen by integrating along line segments across the
finitely many boxes.

For $\kappa>0$, the vertical slope of
$q_\kappa=(1-\kappa)q_0+\kappa v$ is at least $\kappa$ and at most
two.  The horizontal Lipschitz bound is still uniform.  Each column
map is therefore a bi-Lipschitz increasing homeomorphism fixing its
ends; its inverse depends Lipschitz-continuously on the horizontal
coordinate, with a bound depending on $\kappa$.  Patching with the
identity proves the asserted bi-Lipschitz property of $S_\kappa$.
No inverse bound uniform as $\kappa\downarrow0$ is asserted.

Lipschitz composition and the ACL characterization of Sobolev maps
give the first bound of \eqref{hp:flattening-derivative}.  For the
second, work in the assigned coordinates near a retained source
point and use \eqref{hp:anisotropic-correction} along almost every
source coordinate line.  At almost every such point,
\[
 |\partial_i(R\circ f)|
 \le C\lambda|\operatorname{proj}_{e^\perp}\partial_i f|
       +C|e\cdot\partial_i f|
 \le C\lambda|\partial_i f|,
\]
where Lemma~\ref{hp:transported-bundle} supplies the last inequality.
Summing the three coordinate estimates proves the claim, also on
band junctions: it does not require differentiability of $S_0$ at
the target datum.  The horizontal component of $DF_0$ is exactly
$\operatorname{proj}_{e^\perp}Df$.  The projection is injective on
$V(\nu,f(x))$, so this component has rank $\rank Df$ on the
retained data.  Conversely, Lipschitz composition cannot increase
rank here.  To see this without differentiating $S_0$ at the target
point, combine Fr\'echet differentiation of $f$ with approximate
differentiation of $F_0$: the inequality
$|F_0(x+z)-F_0(x)|\le C|Df(x)z|+o(|z|)$ implies
$|DF_0(x)v|\le C|Df(x)v|$ for every $v$, by taking approximate
limits in narrow cones about $v$.  Thus
$\ker Df(x)\subset\ker DF_0(x)$, proving equality of ranks.
Finally the exact identity
\[
 F_\kappa=(1-\kappa)F_0+\kappa f
\]
and $F_0\in W^{1,p}$ prove strong convergence as $\kappa\downarrow0$.

We next identify the exceptional set of values of $S_0$ exactly.
On a band, when $\chi(h)=1$, its image height is
\[
 b_j(h)=s_j(h)-(j-1)L.
\]
Let
\begin{equation}\label{hp:collapsed-values}
 Z_{\mathrm{box}}=
   \bigcup_{j=1}^N\{(h,b_j(h)):\chi(h)=1\},
 \qquad Z=\bigcup_{\mathrm{boxes}}Z_{\mathrm{box}}.
\end{equation}
These are compact sets inside their boxes.  Each $b_j$ is piecewise
affine with slope at most $C\lambda$, so $Z$ has zero volume.
When adjacent bands touch, their displayed image values coincide;
the resulting fibre is the entire contiguous closed interval.
These and only these intervals are the non-singleton fibres.  Indeed,
if $\chi(h)<1$, the whole column is strictly increasing.  If
$\chi(h)=1$, it is strictly increasing outside the closed union of
bands.

For a value outside $Z$ there is exactly one preimage.  Near that
preimage the vertical slope has a positive lower bound: either
$1-\chi$ is bounded below locally, or the point is at positive
distance from all closed bands.  Together with
\eqref{hp:anisotropic-correction}, this gives a local Lipschitz
inverse.  This proves the precise homeomorphic-complement statement
\eqref{hp:flattening-homeomorphic-complement}.  The selected data
lie where $\chi=1$ and in the union of bands, and hence their images
lie in $Z$.

It remains to arrange genuine separated plate patches, rather than
merely a finite union of sheets.  Refine each piecewise affine graph
into its finitely many affine pieces, including their lower-dimensional
strata.  Assign every output datum to one containing affine plane;
if sheets coincide or meet, make one measurable choice.  Use the
finite pushforward of the retained source weight under $F_0$ for
this assignment.  Inner regularity gives disjoint compact subsets
of the assignments with an arbitrarily small further loss.  These
compacts have disjoint ambient neighborhoods.  In each plane, take
a sufficiently fine rectangular grid, choose its cut lines to have
zero assigned measure, and remove thin neighborhoods of the cuts.
The remaining compact pieces lie strictly inside flat rectangles
with margins, and the rectangles and ambient neighborhoods can be
chosen disjoint using the preceding separation.  Pull back the
retained compact subsets through $F_0$ within the already compact
source data.  This gives the final compact $K$ and all the claimed
patch properties.  The initial chart separation prevents saturation
from mixing different rank or range-line assignments.  In particular,
no claim that sheet intersections or projected singular measures
have zero assigned mass is needed.  On the source subset assigned
to a given plate, the normal component of $F_0$ is constant.
Locality of weak derivatives therefore puts the range of $DF_0$
in that plate's tangent plane almost everywhere on the subset.
Together with rank preservation and
\eqref{hp:flattening-derivative}, this also preserves the prescribed
approximate line direction in the rank-one case, with an additional
$C\lambda$ angular error.  Distribute the losses in this
and the preceding selections so that their sum is less than $\eps$.

Lastly, put $U=\Omega\setminus F_0^{-1}(Z)$.  On $U$, $F_0$ is a
Sobolev homeomorphism with positive orientation onto
$\Lambda\setminus Z$.  Apply Lemma~\ref{hp:regularity} on a
countable interior exhaustion.  Its nonregular source points form
a fixed null set $E\subset U$, and the image of its regular
rank-deficient points is null.  This proves the asserted dichotomy
for almost every target point.  For $p\ge3$, the volume property
$N$ makes $F_0(E)$ null as well.  For $2<p<3$, that image is not
assumed null and is retained as the stated exceptional alternative.
\end{proof}

\subsection{Protected endpoint pairs}

The next step selects a pair of target rectangles on opposite sides of
each retained plate.  Their geometric purpose is to give exactly two
contact points for each slope in a central rectangle.  In the convex
quantizer of Section~\ref{sec:quantizers}, these paired contacts produce
a planar output face containing the eventual quantized images of the
retained data.  The carrier must
first make neighborhoods of the endpoint rectangles exact affine or
PL regions; the central data themselves will remain untouched by those
implants.

We continue to measure losses by the finite measure introduced before
Lemma~\ref{hp:transported-bundle}:
\[
 d\mu(x)=\ind_{\{1\leq\rank Df\leq2\}}(x)
             (1+|Df(x)|^p)\,dx.
\]
Only regular points, as in Lemma~\ref{hp:regularity}, enter this measure.
Every compact restriction and every finite measurable partition below may
discard a prescribed arbitrarily small amount of this finite measure.
In particular, a finite sequence of such restrictions can be made with a
prescribed total loss.  We apply Proposition~\ref{hp:flattening}, and write
$F_0=S_0f$ and $F_\kappa=S_\kappa f$ for its maps.

\begin{lemma}[Protected wells]\label{hp:wells}
Let finitely many separated compact pieces of the retained $F_0$-image
lie on flat plates, with relatively compact disjoint ambient chart
neighborhoods.  Let a positive continuous motion scale $\delta$ on
$\Lambda$ be prescribed.  After an arbitrarily small $\mu$-loss, one can
choose, on each plate, an origin $z$, a unit vector $n$ as close as
desired to its normal, a number $d>0$, and nested protected horizontal
rectangles in $n^\perp$ with the following properties.
There is a nonnegative smooth compactly supported function $b_1$ on
$\Lambda$, extended by zero to $\R^3$, such that
\begin{enumerate}[label=\textup{(\roman*)}]
\item $\Delta b_1\leq C$ with an absolute constant, and $b_1$ is as
small as prescribed relative to $\delta^2$.  In particular
$b_1=o(\dist(\cdot,\R^3\setminus\Lambda)^2)$ at the boundary.
\item For every $h$ in a protected horizontal rectangle, put $q=z+h$.
The affine function
\[
 L_h(x)=q\cdot x-\tfrac12|q|^2+d^2
\]
supports $\psi_1(x)=|x|^2/2+b_1(x)$ globally, and its contact set
consists exactly of $q+dn$ and $q-dn$.  In particular,
\begin{equation}\label{hp:paired-dual-value}
 \sup_x\{q\cdot x-\psi_1(x)\}=\tfrac12|q|^2-d^2.
\end{equation}
For a fixed compact rectangle and any fixed neighborhoods of these
two endpoint sheets, the supporting gap is uniformly positive outside
the neighborhoods.
\item The assigned central data $X=F_0(x)$ have
$|\langle X-z,n\rangle|<d/4$.  For almost every assigned $x$ with
respect to $\mu$, both points
\begin{equation}\label{hp:sampled-endpoints}
 Y_\pm(X)=X+n\bigl(\pm d-\langle X-z,n\rangle\bigr)
\end{equation}
belong to genuine homeomorphic target neighborhoods of $F_0$ and have
either a regular full-rank inverse point or an inverse point in a fixed
source null set.  The latter alternative is needed only when $2<p<3$.
\end{enumerate}
All assertions persist after shrinking the protected horizontal
rectangles.  The endpoint assignments can be restricted to compact
sets on which their inverse points and their regular or exceptional
types have fixed separated compact ranges.
\end{lemma}

Figure~\ref{hp:fig-paired-endpoints} shows the geometry of these pairs.
The endpoint rectangles will become exact target regions in the first
round of Proposition~\ref{hp:carrier}; the present lemma selects their
positions and inverse types.

\begin{figure}[htbp]
\centering
\begin{tikzpicture}[x=1.05cm,y=.9cm,>=Latex,font=\small]
  \fill[orange!8] (-2.9,1.05) rectangle (2.9,1.55);
  \fill[orange!8] (-2.9,-1.55) rectangle (2.9,-1.05);
  \draw[orange!65!black,dashed] (-2.9,1.05) rectangle (2.9,1.55);
  \draw[orange!65!black,dashed] (-2.9,-1.55) rectangle (2.9,-1.05);
  \draw[orange!80!black,very thick] (-2.5,1.3)--(2.5,1.3);
  \draw[orange!80!black,very thick] (-2.5,-1.3)--(2.5,-1.3);
  \draw[gray,densely dotted] (-3.1,0)--(3.1,0);
  \draw[blue!65!black,very thick] (-2.6,-.4)--(2.6,.4);
  \draw[gray,dashed] (1.2,-1.3)--(1.2,1.3);
  \fill (0,0) circle (1.4pt) node[below left] {$z$};
  \fill (1.2,0) circle (1.4pt) node[below right] {$q=z+h$};
  \fill[blue!65!black] (1.2,.185) circle (1.8pt)
      node[above left=2pt] {$X=F_0(x)$};
  \fill[orange!80!black] (1.2,1.3) circle (1.8pt)
      node[above=5pt] {$Y_+(X)=q+dn$};
  \fill[orange!80!black] (1.2,-1.3) circle (1.8pt)
      node[below=5pt] {$Y_-(X)=q-dn$};
  \draw[->] (-4.7,-1.5)--(-4.7,1.65) node[above] {$n$};
  \node[anchor=east] at (-3.05,1.3) {$v=d$};
  \node[anchor=east] at (-3.05,-1.3) {$v=-d$};
  \node[blue!65!black,anchor=west] at (2.8,.48) {retained plate};
  \node[anchor=west] at (3.2,0) {$v=0$};
\end{tikzpicture}
\caption{A section through a retained plate and its paired endpoint
rectangles; horizontal rectangles appear as intervals.  The chosen
normal $n$ may be slightly tilted from the plate normal, so the datum
$X=F_0(x)$ need not lie in the central plane $v=0$.  Both endpoints have
the same horizontal coordinate $h$.  The dashed boxes indicate
neighborhoods that the carrier will later make exact.  The diagram
shows $F_0$ data, before the positive parameter $\kappa$ is chosen.}
\label{hp:fig-paired-endpoints}
\end{figure}


\begin{proof}
First work on one plate.  Choose a horizontal cutoff $0\leq\theta\leq1$
which is one on a neighborhood of a compact rectangle strictly larger
than the rectangle to be protected.  Its support is compact in the
horizontal chart.  Choose a smooth nonnegative function $\beta$,
supported in $[-2,2]$, such that
\[
 u\longmapsto u^2/2+\beta(u)
\]
has minimum $1$ exactly at $u=\pm1$.  For example, first prescribe this
function to be $1+(u^2-1)^2$ for $|u|\leq3/2$; in the remaining
transition region the unperturbed function $u^2/2$ is already larger
than $1$, so the nonnegative bump can be smoothly cut off there without
creating another minimum.  In the coordinates $x=z+h+vn$, set
\begin{equation}\label{hp:well-formula}
 b_z(x)=\theta(h)d^2\beta(v/d).
\end{equation}
Choose $d$ much smaller than the horizontal cutoff margin and the
ambient chart margin.  All supports can then be made disjoint.  Since
\[
 \Delta b_z=d^2\beta(v/d)\Delta_h\theta+
                     \theta(h)\beta''(v/d),
\]
choosing $d^2\|\Delta_h\theta\|_\infty\leq1$ gives an absolute upper
bound for $\Delta b_z$.  Disjointness gives the same bound for their
sum $b_1$.  Its size is $O(d^2)$ on each chart, so the finitely many
$d$'s may also enforce every prescribed $\delta^2$ bound.  The support
is compact in $\Lambda$, proving the stated boundary decay.

For a protected $h_0$, subtract $L_{h_0}$ and write $q=z+h_0$.  Within
the chart the difference is
\[
 \psi_1(x)-L_{h_0}(x)
 =\tfrac12|h-h_0|^2+\tfrac12v^2+b_1(x)-d^2.
\]
Where $\theta=1$, the one-dimensional choice of $\beta$ makes this
nonnegative, with equality precisely when $h=h_0$ and $v=\pm d$.
Where $\theta\ne1$, the horizontal margin has been chosen so that
$|h-h_0|^2/2>d^2$.  The other bumps are nonnegative and have disjoint
supports.  The same distance estimate applies outside the chart.
This proves the global contact statement and
\eqref{hp:paired-dual-value}.  On compact $h_0$-sets the gap is
uniform away from the endpoint neighborhoods: otherwise a sequence
of almost contacts would, by quadratic growth and compactness, limit
to a third contact.

Initially the central data lie on the original flat plate.  Choose $d$
in a small positive interval $[d_-,d_+]$, and then restrict $n$ to a
sufficiently small open cap about the original normal that
$|\langle X-z,n\rangle|<d_-/4$ for all central data in the compact
patch.  All preceding supports, cutoffs and margins may be chosen
uniformly over this parameter range.  Give $n$ and $d$ smooth
probability densities in the cap and interval.  For fixed $X$, the map
in \eqref{hp:sampled-endpoints} has the form $X+r n$, where
$r=\pm d-\langle X-z,n\rangle$.  In sphere coordinates its volume
Jacobian has absolute value $r^2$: differentiation in $d$ supplies
$\pm n$, and each angular derivative has tangential part $r\,dn$.
Here $|r|\geq3d_-/4$.  Thus each endpoint law is absolutely continuous
for three-dimensional volume.

By Proposition~\ref{hp:flattening}, almost every target point off the
closed sheet-value set has exactly the stated regular or exceptional
type.  Fubini applied to the finite assignment measure therefore gives
a parameter choice for which both endpoints have an allowed type for
$\mu$-almost every assigned point.  This is an averaging argument in
three target dimensions; it uses no absolute continuity of the
horizontal assignment measure.  Push that measure to the horizontal
patch.  Split according to the two endpoint types simultaneously and
restrict each part to compact subsets of its genuine inverse charts.
The inverse endpoint maps are continuous there.  Further compact
restriction gives separated compact source ranges for the finitely
many classifications, with arbitrarily small loss.  The central
compact data remain separated from all endpoint inverse data because
their $F_0$-values lie on the central sheet, whereas the endpoints lie
in its homeomorphic complement.  A finite rectangular refinement
with smaller margins supplies the rectangles in the statement.
\end{proof}

\subsection{Implanting exact cores by cubical collapse}

The endpoint selection has located two compact sets of target values
over each retained plate.  Their inverse data are either regular full
rank or belong to a source null set, but the maps near these values are
not yet required to be affine or PL.  We now construct such exact
neighborhoods.  The regular data will permit an energy estimate with a
universal constant.  For exceptional inverse data, which occur only
when $2<p<3$, the construction will confine the uncontrolled cost to
marked source cubes.  The data are selected using $F_0$, while the
initial replacements will be made in the homeomorphism $F_\kappa$
for a sufficiently small $\kappa>0$.

The local replacement has a simple normalized goal.  Given a
homeomorphism close to the identity near the boundary of a cube, make
it exactly the identity on the inner cube, leave its outer boundary
values and its image unchanged, and control the energy in the
intervening shell.  The output is allowed to collapse isolated tame
balls.  Those fibres can subsequently be opened while preserving
previously exact target regions; the opening itself supplies no energy
estimate.

The energy estimate comes from sampling the input near a
two-dimensional cubical skeleton.  Small source charts of size $\ell$
are immersed at one common scale $s$; the factors in the derivative of
the lift cancel, and the weighted overlaps of the samples are summable.
We first construct that immersed skeleton, then carry out the
replacement and identify its fibres.

Write
\[
 B=[-1,1]^3,\qquad t(x)=\max_{1\leq j\leq3}|x_j|-1,
 \qquad B_r=\{x:t(x)<r\}.
\]
Thus the notation $B$ denotes the closed inner cube, whereas $B_r$ is
open.  Boundaries have zero volume, so this convention does not affect
any integral below.  The constants in the next two Lemmas depend on the fixed cubical
construction and, when indicated, on $p$; they do not depend on the input
homeomorphism or on the sufficiently small dyadic number $s$.

\begin{lemma}[An immersed cubical skeleton]\label{hp:immersed-skeleton}
There are positive constants $b,C_0,C_1$ and, for every sufficiently small
dyadic $s$, a locally finite conforming triangulation $\mathcal T_s$ of
$\{0<t<2s\}$ with the following properties.  If $D\in\mathcal T_s$ has
scale $\ell_D$, then
\[
 \diam D\asymp\ell_D,
 \qquad \ell_D\asymp\min\{s,t(x)\}\quad(x\in D),
\]
where the second assertion allows the uniform constants implicit in
$\asymp$.  The simplices have uniformly controlled shapes and finitely
many local configurations up to translation and rescaling.  There is a
smooth orientation-preserving local diffeomorphism $i$ on an open
neighborhood $\mathcal U$ of their two-skeleton, including a neighborhood
of the outer boundary $\{t=2s\}$, such that
\[
 i=\Id\text{ near }\{t=2s\},\qquad
 |i(x)-x|\leq C_0s.
\]
The neighborhood contains, at every simplex, a neighborhood of its
skeleton of width $b\ell_D$.  On a slightly smaller such neighborhood,
\begin{equation}\label{hp:immersion-bounds}
 |Di|\leq C_1\frac{s}{\ell_D},\qquad
 |(Di)^{-1}|\leq C_1\frac{\ell_D}{s},\qquad
 |D^2i|\leq C_1\frac{s}{\ell_D^2}.
\end{equation}
There are also uniform local inverse charts: a chart centered at a
point of the smaller neighborhood contains a source ball of radius
$b_1\ell_D$, and its image contains the corresponding target ball of
radius $b_2s$, for fixed $b_1,b_2>0$.

More precisely, after splitting into a bounded number of charts per
simplex, the formulas for $i$ have the form
\begin{equation}\label{hp:immersion-template}
 i(x_*+\ell\xi)=y_*+sI(\xi),
\end{equation}
where $I$ belongs to a fixed finite list of smooth formulas on fixed
compact chart neighborhoods, and $y_*$ belongs to a fixed rational
subdivision of the $s$-grid.  For fixed $s$ there are only finitely many
possible $y_*$ in the bounded region used by this construction.  The
map $i$ is permitted to have overlaps between different local charts.
\end{lemma}

\begin{proof}
Set $\lambda_k=2^{-k}s$, $k\geq0$.  Tile
\[
 \{\lambda_k<t<2\lambda_k\}
\]
by axis-aligned cubes with side comparable to $\lambda_k$; a fixed
dyadic subdivision can be used throughout.  This is possible with
exact alignment because $1$, $s$, and all $\lambda_k$ are dyadic.
Across consecutive layers, a face of a coarser cube is subdivided to
match its finitely many finer neighbors.  Triangulate each resulting
face by coning its subdivided boundary to its center, and cone these
triangles to the center of the cube.  Shared faces receive the same
triangulation.  The neighbor-size ratios are bounded, and all vertex
positions in a normalized cube belong to a fixed finite set.  This
proves the shape and local-configuration assertions.  The harmless
finite modifications needed near $t=2s$ can be made on fixed dyadic
levels.  In particular a thin outer band can be reserved for the
identity map.

We construct the immersion successively near vertices, edges, and
faces.  All neighborhoods in this construction are in the source;
there is no requirement that their images be disjoint.  Choose at each
vertex $v$ a representative incident scale $\ell_v$.  The ratios between
incident scales belong to a finite list.  Away from the reserved outer
band, choose a nearest point $y_v$ of the $s$-grid and prescribe
\[
 i(x)=y_v+\frac{s}{\ell_v}(x-v)
\]
on a small vertex neighborhood.  These source neighborhoods are
disjoint when their relative radii are small enough.  On the outer
band instead prescribe the identity germ.  Its normalized scale ratio
is bounded above and below, since $\ell_v\asymp s$ there.  All target
anchors, including those in the identity band, lie on a fixed rational
subdivision of the $s$-grid.  Adjacent anchors differ by at most $Cs$.

Consider an edge core outside the vertex neighborhoods.  Its two end
germs are already prescribed.  Join these germs by a regular immersed
arc in a ball of radius $Cs$ containing them.  This can be done relative
to the end intervals by the one-dimensional relative immersion
Theorem.  Choose two transverse vector fields along the arc, equal to
the prescribed transverse derivatives on the end intervals.  They
give a positive three-frame.  Relative to its positive scalar end
frames, the homotopy class of this frame path can be made trivial:
one full rotation of the two transverse vectors changes the class in
$\pi_1(GL^+(3))\cong\mathbb Z/2$.  Thus choose the rotation, if needed,
so that the entire frame path is homotopic to the path of positive
scalar frames, with its ends fixed.  Thickening the framed arc gives
an orientation-preserving local diffeomorphism on a sufficiently thin
edge neighborhood.  It agrees exactly with the old affine formulas on
the end collars.  Distinct truncated edge cores have disjoint source
neighborhoods.

For a triangular face, remove the already treated vertex and edge
neighborhoods.  Choose a rounded disk in the remaining face whose
boundary collar lies in those neighborhoods.  The prescribed local
diffeomorphism on that collar provides both a map of the collar and a
positive three-frame.  Its frame loop is null-homotopic.  Indeed, move
the loop within the treated neighborhood onto the triangle perimeter:
each edge frame path has the trivial relative class just prescribed,
and the vertex frames are positive scalars.  Concatenating these paths
therefore gives the trivial loop in $GL^+(3)$.

The target ball is contractible, so the collar map extends as a smooth
base map of the disk into a slightly larger ball of radius $Cs$.
The null-homotopy of the three-frame extends its restriction to the two
fixed tangent directions of the face.  We have consequently obtained
a formal immersion of the disk that is the differential of the
prescribed map near its boundary.  The relative Smale--Hirsch Theorem,
in source dimension two and target dimension three, gives an actual
immersed disk agreeing with the prescribed map on a smaller boundary
collar; see \cite[Theorem~5.7]{Hirsch1959}.  We use the target open ball
as the target manifold in that Theorem.  Thus the extension stays in
a ball of radius $Cs$.  In particular this is not an application of a
relative immersion extension theorem in equal dimensions.

To thicken the disk, choose a positive transverse column $\nu$ along it,
equal to the old normal derivative on the boundary collar.  This
choice extends because the set of positive transverse columns over
an oriented tangent two-frame is a contractible half-space.  In flat
source coordinates $(u,r)$ use the first-order formula
\[
 \phi(u)+r\nu(u).
\]
It is a local diffeomorphism for sufficiently small $|r|$.  On the
boundary collar, its difference from the old exact formula is
$O(r^2)$, with derivative difference $O(|r|)$.  A fixed smooth cutoff
in $u$ blends it with that old formula.  After reducing the normal
width, the differential remains positive and nonsingular, and the
formulas agree exactly on the region retained from the old
neighborhood.  Distinct face cores have disjoint source neighborhoods;
their overlaps with the old neighborhood are precisely the regions
where this matching has been imposed.  This constructs a smooth
positive local diffeomorphism on a joint neighborhood of the entire
two-skeleton.

We now justify all uniformity assertions, which are not consequences
of the qualitative immersion Theorem alone.  Normalize an incident
source configuration by its scale $\ell$ and normalize target lengths
by $s$.  There are finitely many source configurations and incident
scale ratios.  Differences between the target anchors are bounded
integer vectors on the fixed rational grid, so they also have only
finitely many possibilities.  The identity prescriptions occur only
in the outer layers, where the normalized ratios and positions have
finitely many possibilities.  First choose one edge extension for each
of these finitely many edge data, including its collar formula and
framing.  Next choose one face extension for each of the finitely many
face data determined by those choices.  Fix these choices once and
for all.  This gives a finite list of normalized smooth formulas on
compact subneighborhoods.

On each such compact subneighborhood the least singular value is
positive and all first and second derivatives are bounded.  Take the
minimum of the finitely many positive margins and the maximum of the
finitely many derivative bounds.  Shrinking once more gives a common
relative neighborhood width.  The inverse function Theorem, with
these uniform bounds and a common reduction of the chart radius,
gives source inverse-chart radii comparable to $\ell$ and image radii
comparable to $s$.  Scaling back gives
\eqref{hp:immersion-bounds} and
\eqref{hp:immersion-template}.  Finally every formula lies within
$Cs$ of its anchor, and every anchor is within $Cs$ of its source
configuration.  Hence $|i(x)-x|\leq C_0s$.
\end{proof}

\begin{lemma}[Cubical cutoff]\label{hp:cutoff}
There are constants $C>2$ and $c>0$ with the following property.
Let $1\leq p<\infty$, let $s$ be sufficiently small and dyadic, and
let $h$ be a homeomorphism on an open neighborhood of
$\overline{B_{Cs}}$, onto its image, with
$h\in W^{1,p}_{\mathrm{loc}}$.  Suppose that
\begin{equation}\label{hp:cutoff-closeness}
 |h(x)-x|\leq cs\qquad\text{whenever }|t(x)|\leq Cs.
\end{equation}
One can replace $h$ inside $B_{2s}$ by a continuous Sobolev map $F$,
keeping a neighborhood of its boundary unchanged, so that
\begin{align}
 &F=\Id\quad\text{on }B,
 \qquad F(B_{2s})=h(B_{2s}),\label{hp:cutoff-image}\\
 &\int_{0<t<2s}|DF|^p
 \leq C_p\int_{|t|\leq Cs}|Dh|^p.\label{hp:cutoff-energy}
\end{align}
The only nonsingleton fibers of $F$ are tame closed balls over a
countable set of target points.  Their diameters tend to zero at every
accumulation, and those target points can accumulate only on
$\partial B$.  These fibers can be opened simultaneously to give a
homeomorphism, altering $F$ only over arbitrarily small mutually
disjoint target neighborhoods of their values.  This last assertion
is topological; it does not assert a Sobolev energy bound for the
opened maps.  If $h$ has volume Lusin property $N$, then so does $F$.
\end{lemma}

\begin{proof}
\emph{The lift and its energy.}
Apply Lemma~\ref{hp:immersed-skeleton}.  Increase $C$ so that all
auxiliary image neighborhoods and inverse-chart neighborhoods are
contained in $\{|t|<Cs\}$.  This is possible because their centers have
$0\leq t\leq2s$ and their radii are at most a fixed multiple of $s$.
Put $\rho(x)=\min\{s,t(x)\}$ in the open shell.  This is a Lipschitz
scale function and is comparable to $\ell_D$ on every simplex $D$.

On a smaller skeleton neighborhood define $H(x)$ by the near-$x$
inverse branch of
\begin{equation}\label{hp:cutoff-lift}
 i(H(x))=h(i(x)).
\end{equation}
The right side differs from $i(x)$ by at most $cs$.  Uniform inverse
charts in Lemma~\ref{hp:immersed-skeleton} therefore make this
definition possible when $c$ is sufficiently small and give
\begin{equation}\label{hp:lift-displacement}
 |H(x)-x|\leq A c\rho(x).
\end{equation}
The branches agree on overlaps.  To see this, reduce $c$ so that any
two candidate values near a given source point lie in the same
injectivity chart; continuity of the local inverse then gives agreement
on overlapping neighborhoods.  In particular $H$ is continuous and
is locally a composition of homeomorphisms.

It is also globally injective on the smaller skeleton neighborhood.
Suppose $H(x)=H(x')$ and write $\delta=Ac$.  Then
\[
 |x-x'|\leq\delta\bigl(\rho(x)+\rho(x')\bigr).
\]
Since $\rho$ is Lipschitz, choosing $\delta<1/10$ makes $\rho(x)$ and
$\rho(x')$ comparable and gives $|x-x'|\leq3\delta\rho(x)$.
With a further fixed reduction of $c$, the points $x,x'$ and their
common lifted image lie in one injectivity chart for $i$.  Equation
\eqref{hp:cutoff-lift} and injectivity of $h$ first give
$i(x)=i(x')$, and injectivity of that chart then gives $x=x'$.
Thus $H$ is an embedding.  On the reserved outer band $i=\Id$; the
same is true at the nearby lifted points when $c$ is small.  Consequently
$H=h$ on a neighborhood of $\{t=2s\}$.

The weak chain rule in a local chart gives
\[
 DH(x)=Di(H(x))^{-1}\,Dh(i(x))\,Di(x)
 \qquad\text{for almost every }x.
\]
The two scale factors in \eqref{hp:immersion-bounds} cancel, so
$|DH(x)|\leq A|Dh(i(x))|$.  Each used part of a simplex is covered by
a uniformly bounded number of injectivity charts of $i$.  Change
variables separately in those charts.  Their Jacobians satisfy
$|\det Di|\geq a(s/\ell_D)^3$, and their images lie in
$B(x_D,As)$ for any fixed $x_D\in D$, after increasing $A$.  We obtain
\begin{equation}\label{hp:lift-cell-energy}
 \int_{D\cap\operatorname{dom}H}|DH|^p
 \leq A_p\left(\frac{\ell_D}{s}\right)^3
       \int_{B(x_D,As)\cap\{|t|\leq Cs\}}|Dh(y)|^p\,dy.
\end{equation}
No global injectivity of $i$ has been used in this change of variables.

For completeness, fix a dyadic depth $\lambda\leq s$ and a target
point $y$.  Simplices of scale comparable to $\lambda$ whose balls
$B(x_D,As)$ contain $y$ lie in a cubical layer of thickness
$O(\lambda)$ within an $O(s)$-ball.  This set has volume at most
$As^2\lambda$, including near edges and corners of the cube.  Since
the simplex interiors are disjoint and have volume at least
$a\lambda^3$, there are at most $A(s/\lambda)^2$ such simplices.
Summing \eqref{hp:lift-cell-energy}, and then interchanging the
nonnegative sum and integral, gives
\begin{align}
 \sum_D\int_{D\cap\operatorname{dom}H}|DH|^p
 &\leq A_p\sum_{k\geq0}
        \left(\frac{2^{-k}s}{s}\right)^3
        \left(\frac{s}{2^{-k}s}\right)^2
        \int_{|t|\leq Cs}|Dh|^p\notag\\
 &\leq A_p\int_{|t|\leq Cs}|Dh|^p.
 \label{hp:lift-total-energy}
\end{align}
This summation is the quantitative reason for using the immersion:
the local sampling volume contributes a third power, whereas the
number of overlapping samples in one layer contributes only a second
power.

\emph{Topological hole filling.}
We next fill the parts of the simplices where the lift was not
constructed.  These fillings will supply a homeomorphism of the shell;
we will then collapse their images to remove their uncontrolled energy.  In each normalized simplex choose nested smooth convex
balls, obtained by smoothing and slightly shrinking that simplex,
\[
 V_D\Subset V_D^+\Subset U_D\Subset D.
\]
They can be chosen homothetic about a common interior point after
fixing one smooth convex model for each simplex shape.  Choose $V_D$
large enough that $D\setminus\overline{V_D}$ and a two-sided collar
of $\partial V_D$ lie in the lift neighborhood.  The three successive
margins are fixed positive multiples of $\ell_D$.  Taking $c$ small
relative to those margins gives
\[
 H(\partial V_D)\Subset V_D^+.
\]
The sphere $H(\partial V_D)$ is locally flat, since $H$ is an embedding
of a neighborhood of $\partial V_D$.  By the generalized Schoenflies
Theorem, it bounds a tame closed ball $K_D$; see
\cite[Theorem~5]{Brown1960}.  Its bounded side lies in $V_D^+$:
the complement of the convex ball $V_D^+$ is connected to infinity
and misses the sphere.  Extend the boundary homeomorphism
$H|_{\partial V_D}$ to a homeomorphism
$\overline{V_D}\to K_D$.

These extensions fit on the correct sides of the lifted boundaries.
Indeed the open shell with all closed hole interiors removed is
connected through the skeleton to the outer band.  Its lifted image
misses every $H(\partial V_D)$ by injectivity of $H$.  For a fixed
$D$, an outer-band point can be chosen outside $V_D^+$, so this
connected image lies on the exterior side of that sphere.  Different
balls $K_D$ are disjoint, since they lie inside different simplex
interiors.  Therefore filling all holes and using $H$ elsewhere gives
an injective map on the shell.

Extend it by the identity on $B$ and call the resulting map $G$.
This extension is continuous.  On the lift region its displacement
is $O(\rho)$ by \eqref{hp:lift-displacement}; within a filled hole,
both the source and image lie in the same simplex of diameter
$O(\ell_D)=O(t)$.  Thus $G(x)-x\to0$ as $x\to\partial B$ from the
shell.  Lifted points stay outside $B$ because their displacement is
less than $t(x)/2$, and filled points stay outside $B$ because
$K_D\Subset D$.  Hence there is no new failure of injectivity at
the inner boundary.

The map $G$ is consequently a continuous injection on
$\overline{B_{2s}}$, with boundary values $h$.  It is a homeomorphism
onto its compact image.  Invariance of domain and separation by the
boundary sphere show that
\[
 G(B_{2s})=h(B_{2s}):
\]
both are the bounded component enclosed by the same embedded boundary
$h(\partial B_{2s})$.  Thus $G$ glues to $h$ outside this cube.  At
this point $G$ need not have a finite Sobolev energy in the holes.

\emph{Collapsing the uncontrolled fillings.}
We eliminate that uncharged energy by a map in the output coordinates.
Choose the balls $U_D$ so that the ones in the outermost layer leave
a fixed positive multiple of $s$ to $\partial B_{2s}$.  All other
layers are farther from that boundary.  By
\eqref{hp:cutoff-closeness} and a boundary-degree argument,
$B_{2s-cs}\subset h(B_{2s})$.  Reducing $c$ therefore ensures
$U_D\Subset h(B_{2s})$ for every $D$.  The same choice leaves a
whole outer band free of these supports.

On $U_D$ construct a Lipschitz radial squeeze $Q_D$ that is constant
on $\overline{V_D^+}$, equals the identity near $\partial U_D$, and
is a homeomorphism from
$U_D\setminus\overline{V_D^+}$ onto $U_D$ with the collapsed point
removed.  Explicitly, in the uniformly controlled polar coordinates
of the smooth convex model, replace the radial variable by a
continuous function that is zero up to the radius of $V_D^+$,
increases linearly on the next radial interval, and equals the
original radius near $\partial U_D$.  The fixed relative buffer
widths give a Lipschitz constant independent of $D$ and $s$.
Extend by the identity outside $U_D$.  The supports are disjoint.
Their diameters tend to zero toward $\partial B$, so the resulting
map $Q$ is continuous there and equals the identity on $B$.
The uniform Lipschitz bound holds across the seams as well, by
decomposing any line segment into the portions inside the supports
and their complement.  The map $Q$ maps $h(B_{2s})$ onto itself.

Set $F=Q\circ G$.  On a neighborhood of every filled hole this map
is constant, because $K_D\Subset V_D^+$.  Elsewhere it is locally
the Lipschitz composition $Q\circ H$, and
\[
 |DF|\leq\Lip(Q)|DH|\quad\text{almost everywhere there}.
\]
The local Sobolev statements glue across hole boundaries because the
map is constant on two-sided collars of those boundaries.  They glue
across simplex faces because the lift formulas already agree there.
Equation~\eqref{hp:lift-total-energy} proves
\eqref{hp:cutoff-energy} on the open shell.

There is also Sobolev gluing at its limiting inner boundary.  One
direct verification uses absolute continuity on lines.  For almost
every line parallel to a coordinate axis, the restrictions in each
open component of the shell are locally absolutely continuous and
their derivatives are integrable, by Fubini and the energy bound.
An absolutely continuous function on compact subintervals with
integrable derivative and a continuous endpoint limit extends
absolutely continuously to that endpoint.  The line intersects the
cube boundary only finitely many times, except for a null family of
lines contained in boundary planes.  Since $F$ is continuous and is
the identity inside the cube, these restrictions glue to an
absolutely continuous restriction on the whole line.  This proves
the asserted Sobolev regularity and identifies the weak derivatives.
The exterior seam is unchanged on a neighborhood, so no additional
gluing issue occurs there.

\emph{Fibres and relative opening.}
The collapsed fibers have an exact description:
\[
 F^{-1}(a_D)=G^{-1}(\overline{V_D^+}),
\]
where $a_D$ is the center to which $Q_D$ collapses.  They are tame
closed balls, since $G$ is a homeomorphism and the output buffers
are tame balls.  Every other fiber is a singleton.  In any compact
subset away from $\partial B$ there are finitely many simplices;
at the remaining accumulation points their scales tend to zero.
Moreover $\diam G^{-1}(\overline{V_D^+})\leq A\ell_D$.  To verify
this last assertion outside the filled hole, use
\eqref{hp:lift-displacement}: a lifted point mapping into $V_D^+$
lies within $O(\ell_D)$ of $D$ and has comparable Whitney scale.
Inside the hole the assertion is immediate.  This proves the stated
null-size and accumulation properties.

Here is the claimed exact opening procedure.  Around each $a_D$ take
an arbitrarily small target ball $W_D$ with closure in $U_D$, mutually
disjoint from the other chosen balls.  Their radii can be reduced
further at will.  On $Q_D^{-1}(W_D)$ replace the radial collapse by
a strictly increasing radial homeomorphism onto $W_D$, agreeing
with $Q_D$ near its boundary; one can first take a small centered
ball in the polar coordinates if necessary.  Leave $Q_D$ unchanged
outside $Q_D^{-1}(W_D)$.  Composing all these changes with $G$ gives
a homeomorphism and changes $F$ only on $F^{-1}(\bigcup_D W_D)$.
The source diameters at accumulating holes tend to zero, so both
the map and its inverse remain continuous at $\partial B$.
Equivalently, one may use a boundary collar and Schoenflies to fill
each of these smaller neighborhoods.  No regularity of these
topological fillings has been used in the energy estimate.

Finally assume that $h$ has property $N$.  On the lift region,
locally $H=i^{-1}\circ h\circ i$, with smooth diffeomorphic charts
on both sides.  Such compositions inherit property $N$, and so
does their Lipschitz postcomposition by $Q$.  The filled-hole
neighborhoods map to countably many points.  On $B$ the map is the
identity, and the common inner boundary has zero volume and maps
to itself.  These countably many local descriptions prove property
$N$ for $F$.
\end{proof}

\subsection{Opening isolated ball fibres}

We record the precise relative form of the topological opening that is
used both during and after the carrier construction.  A null family
below means that only finitely many members meeting a fixed compact
localization have diameter greater than any prescribed positive
number.

\begin{lemma}[Relative opening]\label{hp:opening}
Let $F:\Omega\to\Lambda$ be a continuous proper onto map obtained by
locally finite image-preserving applications of
Lemma~\ref{hp:cutoff} in genuine homeomorphic neighborhoods.  Assume
that subsequent applications avoid earlier fibre values and their
accumulation sets.  Let $P$ be the countable set of non-singleton
fibre values.  Then its fibres are tame closed balls, form a null
family on compact localizations, and their value accumulations have
singleton fibres.  If $C\subset\Lambda$ is relatively closed and
$C\cap P=\varnothing$, there is a homeomorphism
$H:\Omega\to\Lambda$, as finely close in values to $F$ as prescribed,
such that
\begin{equation}\label{hp:relative-opening-agreement}
 H=F\quad\hbox{on }F^{-1}(C).
\end{equation}
One can require agreement on a neighborhood of every compact target
test already separated from $P$.
\end{lemma}

\begin{proof}
For a single cutoff the assertion about fibres is part of
Lemma~\ref{hp:cutoff}: its nontrivial fibres lie inside separate
Whitney tetrahedra, and their diameters tend to zero at the exact
inner seam.  The corresponding target points accumulate only at that
seam, where the fibre is a singleton.  A later operation in a genuine
homeomorphic neighborhood changes none of those fibres.  Local
finiteness of the operations therefore preserves this description and
the null-family property on every compact localization.  In
particular every $v\in P$ is isolated from the other values in $P$
and from their accumulation set.

Choose mutually disjoint small target balls $B_v$ about these values.
For each fixed $v$ their radii can be made smaller than its distance
from $C$, from the other exceptional values and accumulation seams,
and from all fixed compact tests to be preserved.  These are positive
individual clearances; a common lower bound is neither asserted nor
needed.  Enumerate the values and impose in addition radii tending to
zero, with any prescribed fine value-error bound.  Properness gives
upper semicontinuity of compact fibres: if $V$ is a neighborhood of
$F^{-1}(v)$, a sufficiently small $B_v$ satisfies
$F^{-1}(\overline B_v)\subset V$.  Otherwise points outside $V$
mapping to a sequence converging to $v$ would have a subsequence in a
compact preimage, contradicting the definition of the fibre.

The inverse of $\partial B_v$ is an embedded locally flat sphere:
$F$ is a homeomorphism on a neighborhood of this sphere.  It bounds a
tame ball containing the fibre, either directly by the radial squeeze
description in Lemma~\ref{hp:cutoff}, or by the locally flat
Schoenflies Theorem~\cite[Theorem~5]{Brown1960}.  The preceding
upper semicontinuity places this ball inside a tame neighborhood
compactly contained in $\Omega$.  Replace $F$ on the ball by a ball
homeomorphism extending its boundary parametrization.  To keep an
outer collar fixed, first use two nested target balls and do this
only on the inner one; use the old map on the annular collar.

The replacements have disjoint target images, so the resulting map
is one-to-one and onto.  At an accumulation seam its difference from
$F$ tends to zero, since a change inside $B_v$ has size at most
$\diam B_v$.  It is thus continuous everywhere.  It is a
homeomorphism by invariance of domain, and its image is precisely
$\Lambda$.  Alternatively, properness and the null-family property
also give continuity of the inverse at the singleton seams.  The
balls avoid $C$, which proves
\eqref{hp:relative-opening-agreement}.  Enlarging a compact test to
a small closed neighborhood before choosing the balls gives the last
assertion.
\end{proof}

\subsection{Exceptional boxes and regular affine cores}

We apply the cutoff in two ways, according to the endpoint inverse
classification.  For exceptional data the assignment measure may be
singular, so the first lemma selects source cubes with little lost
assignment mass and small volume and facet energy.  The second lemma
uses regular affine jets to obtain an energy bound with a universal
constant; the accuracy needed to apply it may depend on the fixed jet
bounds.

\begin{lemma}[Exceptional inverse data]\label{hp:exceptional-boxes}
Let $E$ be a compact source null set contained in finitely many genuine
homeomorphic neighborhoods of a continuous $W^{1,p}_{\rm loc}$ map
$F$, where $p>2$.  Let $\sigma$ be a finite measure supported on $E$,
and let $g\in L^1_{\rm loc}$ be nonnegative.  For any $a,b>0$, any
prescribed source and target size bounds, and any closed sets separated
from $E$, one can choose finitely many cubes $U$ with disjoint closures,
avoiding those closed sets and compact in the given neighborhoods,
such that, after a $\sigma$-loss less than $a$, the retained data lie
with positive prescribed relative margins in their interiors and
\begin{equation}\label{hp:exceptional-cube-budget}
 \int_{\bigcup U^+}g< b,
 \qquad
 \sum_U\ell_U\int_{\partial U}^{\rm ext}g\,d\cH^2<b.
\end{equation}
Here $U^+$ are sufficiently small enlarged neighborhoods, and the
superscript $\rm ext$ means that each facet is allowed a fixed small
tangential extension.  The facet positions can be chosen to be
$L^1$ Lebesgue points in their normal variable, and their $F$-images
have zero three-dimensional volume.  The relative insets and the
further margins may be prescribed arbitrarily small before the final
grid phase is selected.
\end{lemma}

\begin{proof}
Take an open neighborhood $V$ of $E$, compact in the stated charts and
disjoint from the forbidden sets, on which $\int_Vg$ is arbitrarily
small.  Such a neighborhood exists by absolute continuity of the
integral.  Choose a cubical grid with small pitch $r$; all cubes and
the small facet extensions that meet $E$ then lie in $V$.  Replace
each grid cube by a concentric inset of side $(1-\tau)r$, with
$\tau>0$ small.  Their closures are disjoint.  Average the phase over
$[0,r)^3$.  For any fixed point, the fraction of phases for which it
lies in a gap or an additional relative margin is $O(\tau)$, with
the additional margin included in this bound.  Integrating against
$\sigma$ controls the expected lost assignment mass, without any
assumption on $\sigma$.

For the facet costs, Fubini in the normal coordinate gives
\[
 \mathbb E\sum_U\ell_U
       \int_{\partial U}^{\rm ext}g\,d\cH^2
       \leq C\int_Vg.
\]
The same estimate holds for the finitely many tangentially extended
facets, with a fixed overlap constant.  Choose $\tau$ so that the
expected lost mass is much smaller than $a$, and choose $V$ so that
the expected facet cost is much smaller than $b$.  Markov's
inequality then leaves a set of phases of positive measure on which
both required bounds hold.  Almost every phase additionally has all
facet coordinates as $L^1$ Lebesgue points of the normal slices of
$g$.  Lemma~\ref{hp:regularity} gives zero image volume for these
generic facets, including their finite tangential extensions.
Intersecting with these full-measure phase conditions costs nothing.
Only finitely many cubes meet the compact set $E$.  A final compact
restriction of the retained assignment set gives strict positive
margins.  Decreasing $r$ beforehand enforces the source sizes and,
by continuity of $F$ on the compact localization, the target sizes.
\end{proof}

\begin{lemma}[Regular affine implantation]\label{hp:regular-implant}
Let $F$ be a positively oriented $W^{1,p}_{\rm loc}$ homeomorphism
on a neighborhood of a compact set of regular full-rank source
points.  Suppose its jets and inverse jets are bounded there and its
Fr\'echet and power-mean differentiation tests hold uniformly below
a fixed radius.  In sufficiently small target grid boxes one can
implant exact positive affine cores, with arbitrarily small relative
grid margins and with disjoint enlarged source comparison boxes.
For a box of target side $R$ using the affine jet $b$ at one of its
assigned inverse points, the new energy in its source support is at
most
\begin{equation}\label{hp:regular-implant-cost}
 C_p\frac{|Db|^pR^3}{\det Db}
 \leq C_p\int_{V_b}|DF|^p,
\end{equation}
where $V_b$ is the source comparison box obtained as the affine
preimage of a buffered target box.  The constant depends only
on $p$ and the fixed cubical cutoff construction, not on the condition
number of $Db$.  The choices may be made relative to previously
fixed disjoint compact sets and exact cores.
\end{lemma}

\begin{proof}
Choose a representative $x_b$ whose image belongs to the pertinent
grid box, and put $b(x)=F(x_b)+DF(x_b)(x-x_b)$.  The source preimage
under $b$ of an enlarged target box lies in a ball of radius $C_bR$
about $x_b$.  The Fr\'echet test, followed by the change of variables
$y=b(x)$, makes
\[
 \sup |F b^{-1}(y)-y|/R
 \quad\hbox{and}\quad
 R^{-3}\int|D(F b^{-1})-I|^p
\]
as small as required.  The accuracy is chosen after the finite jet
and inverse-jet bounds; this is where those bounds affect the
construction.  One may also require
\[
 \int_{V_b}|DF-Db|^p
       \leq 2^{-p-1}|Db|^p |V_b|.
\]
Consequently the normalized derivative integral is $O_p(R^3)$,
whereas $\int_{V_b}|DF|^p$ is bounded below by a fixed multiple of
$|Db|^p|V_b|$.  The volume $|V_b|$ is a fixed shape multiple of
$R^3/\det Db$.

Apply Lemma~\ref{hp:cutoff} to $F b^{-1}$ in the target-normalized
coordinates and then precompose its result with $b$.  It is the
identity on the inner normalized core, hence the resulting source
map equals $b$ on the affine-source core.  Multiplication by $Db$
and the source Jacobian give the first bound in
\eqref{hp:regular-implant-cost}; the preceding lower bound gives
the second.  An inverse condition number is not used in this final
comparison.

For disjointness, inset not only the core but every used enlarged
normalized box from its outer grid box.  Make the value-test error
smaller than the resulting target margin.  Then $F$ maps the
affine-source comparison box inside the associated outer grid box.
Since $F$ is one-to-one on these neighborhoods, source comparison
boxes corresponding to disjoint target boxes are disjoint.  Compact
separation handles different inverse charts and all previously
reserved data.  Uniform tests are obtained whenever needed by
compact restriction in the finite assignment measure before the
mesh is chosen.  This argument does not require the grid centers
themselves to be regular image points.
\end{proof}

\subsection{An exact carrier with prescribed residual volume}

We can now build the carrier.  The first round uses the endpoint
assignments to produce exact neighborhoods of both rectangles in each
pair.  Those neighborhoods and their buffers must be fixed before the
quantizer can specify how small the residual target volume should be.
After these volume thresholds are given, the second round adds exact
cores throughout the remaining target, keeping the first-round regions
and the retained source data fixed.

\begin{proposition}[Two-round exact carrier]\label{hp:carrier}
Prescribe a positive continuous source value tolerance
$a:\Omega\to(0,\infty)$ before applying Proposition~\ref{hp:flattening}.
Use there a target tolerance $a_\Lambda$ satisfying
\[
 a_\Lambda(y)\leq\tfrac14 a(f^{-1}(y)),
\]
and then fix the resulting rank data and endpoint protections of
Lemma~\ref{hp:wells}.  Fix also an enlargement factor $A>1$ for the
marked collars, before assigning their facet budgets, and let $\eps>0$.
After a total $\mu$-loss smaller than $\eps$, a first round of choices
gives $\kappa>0$, a compact retained source set $K$, a continuous
proper onto carrier $F_b^1:\Omega\to\Lambda$, an exact open target
set $O_1$, and compact endpoint rectangles strictly inside $O_1$.
These choices precede any prescribed residual-volume thresholds.
Fix a locally finite target Whitney tiling $\mathcal Q$, and choose
arbitrary numbers $a_Q>0$ for $Q\in\mathcal Q$, after $O_1$ and the
endpoint buffers have been determined.  A second round gives
$F_b:\Omega\to\Lambda$ and $O\supset O_1$ such that:
\begin{enumerate}[label=\textup{(\roman*)}]
\item $F_b$ is continuous, proper and onto, belongs to
$W^{1,p}_{\rm loc}$, and has only the countable tame ball fibres of
Lemma~\ref{hp:opening}.  It is a genuine homeomorphism on target
neighborhoods away from their values and accumulation seams.
For $p\geq3$ it satisfies volume Lusin $N$.
\item $O$ is a locally finite union of exact affine or PL core
images.  Every compact subset of such a core has a neighborhood on
which $F_b$ and its inverse are exactly those of its affine or PL
chart.  In particular no point of $O$ has a nontrivial fibre, and
\begin{equation}\label{hp:residual-volume}
 \cL^3(Q\setminus O)<a_Q\qquad(Q\in\mathcal Q).
\end{equation}
The compact endpoint rectangles retain their fixed buffers in $O_1$.
\item $F_b=F_\kappa$ on a neighborhood of $K$.  It meets the
prescribed fine value tolerance relative to $f$.  The choices of
$\kappa$ and of the retained central margins ensure that
$F_\kappa(K)$ stays in the protected central plate neighborhoods.
\item When $2<p<3$, there is a locally finite family of marked
source cubes $U$ with disjoint closures, away from $K$, such that
\begin{equation}\label{hp:small-marked-mass}
 \int_{\bigcup U}(1+|Df|^p)<\eps.
\end{equation}
For $p\geq3$ this family is empty.  Off the marked cubes the energy
has the universal buffered regional estimate described below.
\item Write $C_U(w)$ for a cube-radius collar of physical width $w$
about $\partial U$.  For the prescribed enlargement factor $A$, the
widths $w_U>0$ can be chosen after both rounds, with disjoint enlarged
marked cubes, so that
\begin{equation}\label{hp:marked-collar-budget}
 \sum_U\frac{\ell_U}{w_U}
             \int_{C_U(Aw_U)}|DF_b|^p<\eps.
\end{equation}
The same bound, multiplied by the $p$th power of the Lipschitz
constant, holds after a uniformly Lipschitz output composition.
The enlarged cubes may still have original $(1+|Df|^p)$-mass less
than $2\eps$.
\end{enumerate}

The regional estimate in \textup{(iv)} has the following meaning.
Suppose a locally finite family of regional tests $A_j$ and two
successive small open buffered enlargements $A_j^{[1]}$ and
$A_j^{[2]}$ are prescribed before the implants are made.  All implants
may be chosen so fine that
\begin{equation}\label{hp:regional-energy}
 \int_{A_j\setminus\bigcup U}|DF_b|^p
       \leq C_p\int_{A_j^{[2]}}|Df|^p
       \qquad\hbox{for every }j.
\end{equation}
Here each enlargement has positive local buffer, the supports and
comparison boxes meeting $A_j$ fit inside its first enlargement at
the corresponding round, and $C_p$ is independent of all selected
jet condition numbers, residual thresholds, and PL reference-map
constants.  In particular the tests may be reserved full-rank cut
shells or ordinary regions outside prescribed interiors.  The same
construction works for summably weighted such tests toward the ends
of the open domains.
\end{proposition}

\begin{proof}
We give the choices in order.  The initial flattening tolerance gives
\[
 |F_\kappa(x)-f(x)|
   =(1-\kappa)|S_0(f(x))-f(x)|<\tfrac14a(x)
       \qquad(0\leq\kappa\leq1).
\]
All subsequent replacements keep their old boundary values and their
old local images.  Choose their compact supports so small that the
oscillation of the map being replaced gives
\[
 |F_b^1-F_\kappa|<a/4,
 \qquad |F_b-F_b^1|<a/4.
\]
The positive minorants in Lemma~\ref{pre:local-tolerances} permit these
choices simultaneously on each locally finite family.  Thus the final
carrier satisfies $|F_b-f|<a$.  These choices control the additional
implantation motion; the first displayed bound has already reserved
the flattening motion before the rank and endpoint data were fixed.

\emph{First round: fixing the exceptional boxes.}
Classify both endpoint signs simultaneously as in
Lemma~\ref{hp:wells}.  In the range $2<p<3$, the exceptional inverse
assignments have compact ranges in a fixed source null set.  Apply
Lemma~\ref{hp:exceptional-boxes}, with $g=1+|Df|^p$, to enclose them
in finitely many marked cubes $U$, discarding only a prescribed
small assignment weight.  Make all these cubes and their enlarged
neighborhoods disjoint from $K$ and from the retained regular inverse
assignments.  Give their source integrals and extended-facet budgets
arbitrarily small shares of $\eps$.  Retain the exceptional inverse
points with strict inner margins.  In the grid phase selection impose
the generic-facet condition for both $F_0$ and the original map $f$.
Lemma~\ref{hp:regularity} supplies full-measure sets of phases for
both conditions, so their intersection retains the same budget choices.
The inset fractions can be as small as desired and are fixed at this
stage.

Inside each $U$, fix a slightly smaller cubical core $U^-$ containing
its retained inverse data with room, and leave enough space for the
cutoff support to remain compact in $U$.  The associated compact
endpoint set $L_U$ lies in $F_0(\interior U^-)$.  Choose a compact
target neighborhood $W_U$ of $L_U$ inside this open image and inside
the genuine homeomorphic target neighborhood of $F_0$.  These margins
are chosen using $F_0$, before $\kappa$ or a PL reference map is fixed.

\emph{First round: fixing the regular grids and all margins.}
First compact-restrict the regular endpoint assignments so that their
jets, inverse jets and differentiation radii for $F_0$ are uniform.
Choose small relative grid margins.  On each plate we will use the
same shifted horizontal square grid for both endpoint signs.  Choose
its pitch $R$ sufficiently small that every pertinent square has a
fixed classification for each sign, using compact separation of the
classification sets.  For an exceptional sign, require that the whole
inset endpoint rectangle lies in the interior of its previously
chosen $W_U$.  For regular signs, the uniform differentiation tests
in Lemma~\ref{hp:regular-implant} permit the required normalized value
and gradient accuracy on every buffered box.  These upper bounds on
$R$ are independent of the grid phase.

Now select that phase.  For every fixed datum, the fraction of phases
lost to the grid margins is bounded by the relative margin fraction.
Fubini therefore gives a phase losing arbitrarily little of the finite
horizontal assignment measure.  Trim compactly inside its surviving
squares, and protect smaller rectangles once more for the actual
central data.

For each square with a regular sign, choose a representative endpoint
inverse and its affine jet $b$.  Form the target three-box of horizontal
and normal side $R$, centered at the grid center at height $\pm d$.
The uniform $F_0$ tests make $F_0b^{-1}$ close to the identity on
every buffered box, including in normalized gradient power mean.
Fix these finitely many tests and their affine-source comparison
boxes.  The source boxes are disjoint because their $F_0$-images stay
strictly inside disjoint outer grid boxes.  Prior compact separation
also makes them disjoint from the marked cubes and from $K$.
All geometric tests and central margins are now fixed.

\emph{First round: choosing $\kappa$ and making the implants.}
Take $\kappa>0$ small enough that all the regular value and gradient
tests also hold for $F_\kappa b^{-1}$.  This follows from
$F_\kappa\to F_0$ in values and in $W^{1,p}$ on each fixed compact.
Require also that $F_\kappa(K)$ remains within the nested central
margins and that $F_\kappa^{-1}(W_U)\Subset\interior U^-$ for each
exceptional core.

The latter inverse requirement follows from uniform inverse convergence
on compact target sets in the homeomorphic complement of $F_0$.
Indeed, if inverse points for $F_\kappa$ failed to converge uniformly
to those for $F_0$, properness and the fine motion bound would give a
convergent subsequence limiting to a different preimage under $F_0$.
Uniqueness in the genuine inverse neighborhood excludes this.
The same compact inverse argument applies to sufficiently fine
homeomorphic approximations of the now fixed $F_\kappa$.

Apply Lemma~\ref{sm:relative-pl} to choose a PL homeomorphism
$h_{\rm ref}$ finely approximating $F_\kappa$.  Choose it so fine
that $h_{\rm ref}^{-1}F_\kappa$ satisfies the normalized identity
tolerance of Lemma~\ref{hp:cutoff} on each exceptional cube, and so
that $h_{\rm ref}^{-1}(W_U)\Subset\interior U^-$.
Inverse continuity on the compact localizations supplies the first
test, and the preceding inverse-margin argument supplies the second.
Apply the cutoff with core $U^-$ and support compact in $U$, and
postcompose by $h_{\rm ref}$.  The result equals $h_{\rm ref}$ on
$U^-$, so its exact core image contains the whole protected endpoint
rectangle with room.  Since $h_{\rm ref}$ and its inverse are locally
bi-Lipschitz, this replacement has finite local $W^{1,p}$ energy;
its possibly large constant is used only inside marked cubes.

For each regular sign, use the cutoff construction of
Lemma~\ref{hp:regular-implant} with the fixed jet and the
$F_\kappa$ tests.  The resulting map is exactly $b$
on its core, whose image contains the inset endpoint rectangle.
These supports are disjoint from the exceptional changes, and all
supports avoid a neighborhood of $K$.  The regular cost is
\[
 C_p |Db|^pR^3/\det Db
       \leq C_p\int_{V_b}|DF_0|^p
       \leq C_p\int_{V_b}|Df|^p,
\]
where the final absolute factor is supplied by
Proposition~\ref{hp:flattening}.  The comparison uses the fixed
$F_0$ jet tests; the $F_\kappa$ tests were required to be equally
accurate before applying the cutoff.  This normalization is why no
condition number enters the final energy charge.

Call the first-round map $F_b^1$, and let $O_1$ be the union of
slightly smaller exact core images.  Each is a polyhedral affine or
PL image and has an exact neighborhood with buffer.  The compact
endpoint rectangles lie strictly inside this open set.  The map is
proper, onto and continuous because the finite changes preserve
local images and agree with the old map near their outer boundaries.
Lemma~\ref{hp:cutoff} supplies its local Sobolev regularity and its
ball-fibre description.  It also preserves volume $N$ when $p\geq3$.
The fibre values and their accumulation sets have zero target
volume: the former are countable, and the latter lie on finitely
many exact affine or PL core boundaries.

\emph{Second round: choosing the sets to fill.}
The first-round choices, including $O_1$ and its endpoint buffers,
are now fixed.  Prescribe all the $a_Q$.  In the interior of each
Whitney cube $Q$, remove $O_1$ from the work region and avoid its
boundary, the first-round fibre values and their accumulations, the
images of marked-cube boundaries, and $F_b^1(K)=F_\kappa(K)$.
Except for $O_1$, all these exclusions have zero volume.  The last
one is null because the retained rank-deficient regular image under
$f$ is null and $S_\kappa$ is Lipschitz.  The first marked-cube boundaries were selected to have null images
under $f$; near them $F_b^1=F_\kappa=S_\kappa f$, and the
Lipschitz map $S_\kappa$ preserves that nullity.

On the remaining genuine homeomorphic portions,
Lemma~\ref{hp:regularity} gives the following exhaustive alternative
at almost every target point: its preimage is regular full rank, or
belongs to a fixed source null set.  Regular deficient images have
zero volume.  When $p\geq3$, volume $N$ eliminates the second
alternative.  Compactly restrict these measurable target sets,
separating their classifications, their inverse charts, and the
inside and outside of the first marked cubes.  The target volume
lost in $Q$ can be made smaller than any fixed fraction of $a_Q$.
All selected sets are compact in the remaining work region.  In
particular their source preimages are compact and all necessary
separations are positive.

\emph{Second round: implants.}
For exceptional assignments outside the old marked cubes use
Lemma~\ref{hp:exceptional-boxes} with
$g=1+|Df|^p+|DF_b^1|^p$.  This is locally integrable, and the
exceptional source set is null.  Choose new marked cubes with
summably small source and facet budgets, disjoint from all earlier
marked cubes, from the regular assignments, and from $K$.  Their
supports have images compact inside the current Whitney interior.
For exceptional assignments already inside an old marked cube,
use smaller reference-map implants there; no new marked cube is
required.  A suitable PL reference is available on the genuine
homeomorphic neighborhoods in use.  One way to obtain it globally
is to apply Lemma~\ref{hp:opening} to $F_b^1$ with a closed
neighborhood of all chosen compact tests fixed, and then use
Lemma~\ref{sm:relative-pl}.  The resulting reference approximates
$F_b^1$ arbitrarily well on these tests.  The same inverse-margin
argument as in the first round makes the chosen endpoint or volume
assignments fall within its exact core images.

For regular assignments use ordinary three-dimensional target grids
and Lemma~\ref{hp:regular-implant}, now with the jets of $F_b^1$.
Compact restriction gives the uniform tests required for a sufficiently
fine grid.  Predetermine arbitrarily small insets, then make the
value and gradient tests correspondingly fine.  Each comparison box
meeting an old marked cube is chosen wholly inside it; exterior
comparison boxes miss every marked cube.  The previously compactly
separated exceptional assignments and new marked cubes are avoided.
The regular source comparison boxes in this round are disjoint, and
their costs are bounded by $C_p\int|DF_b^1|^p$ on those boxes.

For exceptional assignments, grid phase averaging against ordinary
target volume pulled back to the null source set controls the loss
to source insets.  For regular assignments, the target grid margins
have arbitrarily small volume.  Therefore in each $Q$ the retained
assignments can be covered by exact cores with total additional
loss less than the unused share of $a_Q$.  Only finitely many compact
assignments and implants are needed in each Whitney cube.  Whitney
local finiteness and properness then make all second-round supports
locally finite in target and source.  These arguments prove
\eqref{hp:residual-volume}.  They also show that second-round
supports miss a neighborhood of every old marked-cube boundary and
of the compact retained set $K$.

Let $F_b$ be the resulting map and $O$ the union of the old and new
buffered exact core images.  The first-round endpoint cores have
been kept fixed.  Image preservation and local finiteness give
properness and surjectivity.  The local Sobolev, Lusin and fibre
properties follow as in the first round.  The additional
accumulation seams are only the new exact core boundaries, a locally
finite family of zero-volume polyhedra.  Thus
Lemma~\ref{hp:opening} applies to the combined carrier.

\emph{Universal regional charging.}
Outside the marked cubes no reference-map constant enters the
construction.  At the first round disjoint regular comparison boxes
give the energy bound by $C_p\int|Df|^p$; elsewhere
$|DF_\kappa|\leq C|Df|$.  At the second round disjoint regular
comparison boxes give the bound by $C_p\int|DF_b^1|^p$.
If a support or comparison box meets $A_j$, choose it small enough
to lie in $A_j^{[1]}$ at that round.  At the preceding round use
$A_j^{[2]}$ for this first enlargement.  Combining the two disjoint
sum estimates proves \eqref{hp:regional-energy}.  There are only two
rounds, so the number of buffer enlargements and the multiplication
of universal constants are fixed.  Local finiteness of the tests
allows all required support sizes to be chosen simultaneously from
positive local minorants.  For countably many prescribed weighted
tests, take the usual compact exhaustion and summable budgets.
This uses no positive lower bound for a buffer width at the ends
of the domains.

\emph{Marked-cube collars.}
The enlargement factor $A$ was fixed before the marked cubes were
selected.  Choose all initial and added marked-cube budgets so small
that \eqref{hp:small-marked-mass} holds and their extended-facet budgets,
including the factor $C_A$ in the slice estimate below, sum to an
arbitrarily small number.  Near the boundary of a first
marked cube the carrier still equals $F_\kappa$, whose derivative
is bounded by $C|Df|$.  Near a newly marked cube it equals the
unchanged $F_b^1$.  The modifications have positive distance from
each fixed boundary: later supports are locally finite over its
compact image and avoid it.  Thus a possibly very thin unchanged
neighborhood of each $\partial U$ remains.

Choose $w_U$ within that neighborhood.  The collar is contained in
six normal slabs of thickness $O(Aw_U)$ over the extended facets.
The $L^1$ Lebesgue property of the selected normal slices yields
\[
 \limsup_{w\downarrow0}\frac{\ell_U}{w}
       \int_{C_U(Aw)}|DF_b|^p
 \leq C_A\ell_U\int_{\partial U}^{\rm ext}g_U\,d\cH^2,
\]
where $g_U=C|Df|^p$ for old cubes and
$g_U=|DF_b^1|^p$ for added cubes.  Give each cube a further summable
error share and decrease its width to realize this estimate.  The
earlier facet budgets imply \eqref{hp:marked-collar-budget}.
Widths can simultaneously make enlarged cubes disjoint and their
additional original energy and volume arbitrarily small.  Finally,
the Sobolev chain inequality
$|D(L\circ F_b)|\leq\Lip(L)|DF_b|$ proves the assertion about
Lipschitz output compositions.  This completes all parts of the
Proposition.
\end{proof}

\subsection{Full-rank regions reserved for the final correction}

The carrier construction is now established.  The next three results
prepare a separate step in the eventual completion: reserving finitely many full-rank cubes, then making a
preliminary locally bi-Lipschitz approximant affine on those cubes.
In Section~\ref{sec:high-assembly}, the cubes and their energy tests
will be reserved before applying the carrier construction; the actual
correction will be performed only at the end.  Here the replacement must remain a homeomorphism.  We therefore return
to the normalized cube $B=[-1,1]^3$ and the shell coordinate
$t(x)=\max_j|x_j|-1$ of Lemma~\ref{hp:cutoff}, and first prove its
smooth-input version.

\begin{lemma}[Smooth cubical cutoff]\label{hp:smooth-cutoff}
Under the hypotheses of Lemma~\ref{hp:cutoff}, suppose in addition
that $h$ and its inverse are smooth on the neighborhoods in use.
For every $\varepsilon_0>0$ there is instead a locally bi-Lipschitz
homeomorphism $F$ onto the same image, equal to the identity on $B$
and equal to $h$ near and outside $\partial B_{2s}$, such that
\begin{equation}\label{hp:smooth-cutoff-energy}
 \int_{0<t<2s}|DF|^p
 \leq C_p\int_{|t|\leq Cs}|Dh|^p+\varepsilon_0.
\end{equation}
The local bi-Lipschitz constants may depend on $h,s$, and
$\varepsilon_0$.  In particular no uniform inverse Lipschitz bound
is asserted as $\varepsilon_0\downarrow0$.
\end{lemma}

\begin{proof}
Use the construction in Lemma~\ref{hp:cutoff} up to the hole-filling
stage.  The lift is now smooth on its neighborhood, and the lifted
hole boundaries are smooth embedded spheres.  The smooth
three-dimensional Schoenflies Theorem identifies their bounded sides
with smooth balls; see \cite[Section~1.1, Theorem~1.1]{Hatcher2023}.
A prescribed smooth orientation-preserving
sphere parametrization extends over a ball: after a ball
identification, isotope the resulting sphere diffeomorphism to a
rotation, and insert that isotopy in an annulus before using the
rotation in the central ball.  Smale's theorem supplies a jointly
smooth isotopy for the fixed sphere diffeomorphism
\cite[Theorem~6]{Smale1959}; for the stronger smooth dependence on
the sphere diffeomorphism, see also Li--Watts
\cite[Theorem~1.5]{LiWatts2011}.
Using collar coordinates first makes
the extension agree with the already prescribed lift on a smaller
two-sided boundary collar.  Hence each hole has a smooth, and in
particular bi-Lipschitz, filling.  It remains essential to prove that
their constants can be chosen uniformly over the infinitely many
Whitney layers.

Fix $h$ and $s$.  By \eqref{hp:immersion-template}, the normalized
lift in a local chart is
\begin{equation}\label{hp:normalized-smooth-lift}
 \xi\longmapsto
 I^{-1}\!\left(
   \frac{h(y_*+sI(\xi))-y_*}{s}
 \right),
\end{equation}
where the inverse is the specified nearby branch.  Overlaps with
neighboring charts are part of the same finite pattern data.  There
are finitely many choices of $I$ and finitely many possible target
anchors $y_*$, since $s$ is fixed and the sampling region is bounded.
The normalized hole and buffer shapes also have finitely many
choices.  Thus there are only finitely many normalized smooth collar
maps on hole boundaries, not merely a bounded family of arbitrary
boundary maps.  Choose a smooth filling once for each of these
finitely many data.  On their compact closures take the maximum of
the finitely many forward and inverse Lipschitz constants.  Call it
$L(h,s)$.  The physical source and output coordinates of a hole are
both rescaled by its scale $\ell_D$, so this constant does not change
with the depth of the Whitney layer.

The resulting pre-squeeze homeomorphism $G$ therefore has a uniform
finite bound for $|DG|$ in all filled holes, and the analogous
inverse bounds there.  Its lift pieces have such bounds as well:
there are only the finitely many normalized maps
\eqref{hp:normalized-smooth-lift}.  The agreement on collars provides
the same local bounds across all finite interfaces.

Replace the nonstrict radial squeezes by strictly increasing ones.
For a common $0<\alpha<1$, let $Q_{D,\alpha}$ be the homothety of
ratio $\alpha$ about $a_D$ on $V_D^+$, interpolate monotonically to
the identity in the surrounding fixed buffer, and keep the identity
near $\partial U_D$.  Their forward Lipschitz constants are uniformly
bounded independently of $\alpha$, and their inverse constants are
finite for every fixed $\alpha>0$.  On each filled hole,
\[
 D(Q_{D,\alpha}\circ G)=\alpha DG.
\]
Consequently
\[
 \sum_D\int_{V_D}|D(Q_{D,\alpha}\circ G)|^p
 \leq\alpha^pL(h,s)^p\sum_D|V_D|
 \leq\alpha^pL(h,s)^p|B_{2s}\setminus B|.
\]
Choose $\alpha$ so that this is below $\varepsilon_0$.  On the
remaining regions the proof of
\eqref{hp:lift-total-energy} applies without change, with the same
uniform upper bound for the Lipschitz constant of the squeezes.
This proves \eqref{hp:smooth-cutoff-energy}.

For clarity, local bi-Lipschitz regularity includes the inner seam.
For these fixed data the maps and their inverses have uniform finite
local Lipschitz bounds off $\partial B$, by the finite normalized
descriptions and the positive squeeze parameter.  They extend
continuously as the identity on $\partial B$.  On a small line
segment crossing a face, edge, or vertex of that cube, decompose the
segment into its open subintervals off the cube boundary.  Integrate
the common Lipschitz bound on those intervals and use continuity at
their endpoints; on portions inside the cube the map is the identity.
This gives the same finite Lipschitz bound across the seam.  Apply
the identical argument in the image to the inverse, which is also
the identity on $B$.  Thus the final map is locally bi-Lipschitz
throughout the open replacement domain and glues to $h$ as claimed.
\end{proof}

\begin{remark}\label{hp:cutoff-affine-scaling}
The cutoff Lemmas apply in affine output coordinates.  More
explicitly, let
\[
 b(x)=y_0+A(x-x_0),\qquad A\in GL^+(3),\qquad r>0,
\]
and normalize a map $k$ by
\[
 \widehat k(\xi)
 =\frac1r A^{-1}\bigl(k(x_0+r\xi)-y_0\bigr).
\]
If $\widehat k$ satisfies the hypotheses of either cutoff Lemma,
transforming the resulting map back makes it equal to $b$ on
$x_0+rB$ and preserves its old image and its outer boundary collar.
Writing
\[
 S=\{x_0+r\xi:0<t(\xi)<2s\},\qquad
 A^{\mathrm{band}}=\{x_0+r\xi:|t(\xi)|\leq Cs\},
\]
the smooth estimate becomes
\begin{equation}\label{hp:affine-cutoff-energy}
 \int_S|Dk_{\mathrm{new}}|^p
 \leq C_p|A|^p|A^{-1}|^p
          \int_{A^{\mathrm{band}}}|Dk|^p
       +|A|^pr^3\varepsilon_0.
\end{equation}
This follows from $D_\xi\widehat k=A^{-1}D_xk$ and the source
Jacobian $r^3$.  The last additive error can therefore be prescribed
arbitrarily in the original coordinates as well.
\end{remark}

\begin{proposition}[Postponed full-rank correction]
\label{hp:full-rank-correction}
Let $p>2$ and let $f:\Omega\to\Lambda$ be an orientation-preserving
Sobolev homeomorphism between bounded domains.  Let $R$ be the set
of its regular differentiation points at which $Df$ has rank three.
Let $P\Subset\Omega$ be any previously protected compact set
disjoint from $R$.  Given $\varepsilon>0$, one can choose finitely
many pairwise disjoint buffered source cubes, disjoint from $P$,
with centers $x_j\in R$, radii $r_j$, and positive affine jets
\[
 A_j=Df(x_j),\qquad
 b_j(x)=f(x_j)+A_j(x-x_j),
\]
such that the inner cubes $K_j=x_j+r_jB$ satisfy
\begin{equation}\label{hp:full-rank-core-tests}
 \int_{R\setminus\bigcup_jK_j}|Df|^p<\varepsilon,
 \qquad
 \sum_j\int_{K_j}|Df-A_j|^p<\varepsilon.
\end{equation}
Their buffered cutoff bands
\[
 E_j=\{x_j+r_j\xi:|t(\xi)|\leq Cs\}
\]
can have arbitrarily small total $f$-energy, even after multiplication
by all the finitely bounded affine factors
$C_p|A_j|^p|A_j^{-1}|^p$.  On those bands the normalized $f$ is
within $cs/4$ of the identity.

Write $Q_j=x_j+r_jB_{(C+1)s}$ for the buffered outer cubes.
After these choices there are numbers $v_j>0$ and $\beta>0$ such
that every smooth diffeomorphism $k:\Omega\to\Lambda$ satisfying
\[
 \sup_{Q_j}|k-f|<v_j\quad\text{for each }j,
 \qquad \sum_j\int_{E_j}|Dk|^p<\beta
\]
can be changed only inside the associated outer cubes
$x_j+r_jB_{2s}$ to a locally bi-Lipschitz homeomorphism
$\widetilde k:\Omega\to\Lambda$ with
\[
 \widetilde k=b_j\quad\text{on }K_j.
\]
The sum of the derivative-error integrals on the inner cubes and
their modified shells is less than $\varepsilon$, after starting the
selection with smaller error prescriptions if necessary.  The cubes
and tests may also be chosen to retain any prescribed positive
continuous value tolerance for the final modification.

A locally bi-Lipschitz preliminary map in $W^{1,p}(\Omega;\R^3)$
can be used in place of $k$ by first applying Theorem~\ref{sm:smoothing}, preserving the finite
value and band-energy tests with room to spare.
\end{proposition}

\begin{proof}
We spell out the order of choices.  Since $|Df|^p$ is integrable,
choose a compact set $K\subset R\setminus P$, with arbitrarily
small discarded full-rank energy, on which
\[
 |Df(x)|+|Df(x)^{-1}|\leq M
\]
for some finite $M$.  This is possible by inner regularity and
truncation of the finite values of the two matrices.  All points of
$K$ may be taken to be both Fr\'echet differentiation points and
$L^p$ differentiation points of $Df$.  The positive determinant of
their derivatives follows from the orientation and the local-degree
test at an invertible Fr\'echet differential.  Fix $M$ before making
any shell-energy choices.

Take the common dyadic $s>0$ small.  At every $x\in K$ there are
arbitrarily small radii $r$ such that the buffered outer cube
\[
 Q=x+rB_{(C+1)s}
\]
has closure in $\Omega\setminus P$, and, with $A=Df(x)$ and
$b(z)=f(x)+A(z-x)$,
\begin{align}
 &\sup_{z\in Q}|A^{-1}(f(z)-b(z))|
       <\frac{cs}{4}\,r,\label{hp:jet-value-test}\\
 &\int_Q|Df-A|^p<\eta|Q|.\label{hp:jet-energy-test}
\end{align}
Here $\eta>0$ can be prescribed after $M$ and $s$.
The first inequality follows from Fr\'echet differentiability,
using $|A^{-1}|\leq M$; the second follows from $L^p$
differentiation.  Both persist at all sufficiently small radii, so
these outer cubes form a fine Vitali family centered on $K$.
Their radii can additionally be bounded by any prescribed local
upper bounds needed for value accuracy and clearance.

Use Vitali selection and then a finite truncation to obtain disjoint
outer cubes $Q_j$ missing only an arbitrarily small amount of the
weighted measure $\mathbf1_K|Df|^p\,dx$.  The closures may be made
disjoint as well: after the finite selection, shrink the cubes by
arbitrarily small amounts, retaining the strict differentiation
tests and allowing an arbitrarily small additional weighted loss.
Start those tests with strict margins before this shrinking.  The
fixed ratio between an outer cube and its inner cube is preserved
by rescaling its $r_j$.  Set $A_j=Df(x_j)$ and use the notation in
the statement.

Each buffered band $E_j$ and each outer-minus-inner region
$Q_j\setminus K_j$ has volume at most $A s|Q_j|$, with $A$ depending
only on the fixed constant $C$.  Equation
\eqref{hp:jet-energy-test} and
$|X+Y|^p\leq2^{p-1}(|X|^p+|Y|^p)$ give
\begin{equation}\label{hp:full-rank-band-energy}
 \sum_j\int_{E_j\cup(Q_j\setminus K_j)}|Df|^p
 \leq A_p(M^ps+\eta)\sum_j|Q_j|
 \leq A_p(M^ps+\eta)|\Omega|.
\end{equation}
The same differentiation test gives
\begin{equation}\label{hp:full-rank-core-error}
 \sum_j\int_{K_j}|Df-A_j|^p\leq\eta|\Omega|.
\end{equation}
Thus the discarded full-rank energy consists of the initial compact
truncation, the finite Vitali truncation, and a part bounded by
\eqref{hp:full-rank-band-energy}.  This proves
\eqref{hp:full-rank-core-tests} after decreasing all the error
prescriptions.  It also proves the stronger band statement: every
affine cutoff multiplier is at most $C_pM^{2p}$, so choose $s$ and
then $\eta$ to make
\[
 C_pM^{2p}\,A_p(M^ps+\eta)|\Omega|
\]
as small as required.  Notice that $M$ was fixed before $s$ and
$\eta$; there is no bound being chosen after the bands to which it
must apply.  Equation~\eqref{hp:jet-value-test} is exactly the
stated normalized value test for $f$.

Now fix this finite family.  Require the preliminary smooth map
$k$ to satisfy
\begin{equation}\label{hp:preliminary-value-test}
 \sup_{Q_j}|k-f|<\frac{csr_j}{4M}
 \quad\text{for each }j.
\end{equation}
Combining this with \eqref{hp:jet-value-test} shows that
\[
 \widehat k_j(\xi)
 =r_j^{-1}A_j^{-1}\bigl(k(x_j+r_j\xi)-f(x_j)\bigr)
\]
is within $cs/2$ of $\xi$ on $|t(\xi)|\leq Cs$.
In particular it meets the closeness hypothesis of
Lemma~\ref{hp:smooth-cutoff} with a fixed positive margin.  Apply
that Lemma separately in these disjoint cubes, and return to the
original coordinates as in
Remark~\ref{hp:cutoff-affine-scaling}.  Choose the finitely many
additive errors to have arbitrarily small sum.  The new maps are
exactly $b_j$ on $K_j$, agree with $k$ on outer collars, and have
the same images as $k$ in their respective replacement cubes.
Therefore they glue to a locally bi-Lipschitz homeomorphism onto
the same fixed target $\Lambda$.

Let $S_j=\{x_j+r_j\xi:0<t(\xi)<2s\}$ denote the modified shells.
Equation~\eqref{hp:affine-cutoff-energy} yields
\[
 \sum_j\int_{S_j}|D\widetilde k|^p
 \leq C_pM^{2p}\sum_j\int_{E_j}|Dk|^p+\delta,
\]
where $\delta>0$ is freely prescribed.  Hence
\begin{equation}\label{hp:correction-shell-error}
 \sum_j\int_{S_j}|D\widetilde k-Df|^p
 \leq2^{p-1}\left(
 C_pM^{2p}\sum_j\int_{E_j}|Dk|^p+\delta
 +\sum_j\int_{S_j}|Df|^p\right).
\end{equation}
The last term was made small in
\eqref{hp:full-rank-band-energy}.  Choose a positive upper bound
for $\sum_j\int_{E_j}|Dk|^p$ after the finite affine bounds are
known, and take $\delta$ small.  Equations
\eqref{hp:full-rank-core-error} and
\eqref{hp:correction-shell-error} then give the asserted total
derivative accuracy on the corrected regions.

To retain a prescribed positive continuous value tolerance, first
make the cubes small enough that the oscillation of $f$ on each
buffered cube is much smaller than that tolerance there.  This is
allowed by the fineness of the Vitali family and continuity of $f$.
Also impose correspondingly small value errors on $k-f$.
Every modified value lies in the old image of its replacement cube.
Thus its distance from $f$ at the source point is bounded by that
cube's oscillation of $f$ plus the chosen value error of $k$.
These quantities can be made smaller than the requested tolerance.

Finally suppose the available preliminary map is only locally
bi-Lipschitz.  Apply Theorem~\ref{sm:smoothing} with sufficiently
small strong $W^{1,p}$ error and its fine value control on the
finitely many compact buffered cubes.  Strong convergence preserves
their energy upper tests after allowing a strict margin, since
\[
 \int_{E_j}|Dk|^p
 \leq2^{p-1}\int_{E_j}|Dg|^p
      +2^{p-1}\|Dk-Dg\|_{L^p(\Omega)}^p
\]
for a preliminary map $g$.  The fine value control preserves
\eqref{hp:preliminary-value-test}.  The preceding smooth correction
therefore applies.  Any change outside the selected cubes at this
preliminary smoothing stage has the arbitrarily small strong error
prescribed in Theorem~\ref{sm:smoothing}.
\end{proof}

\subsection{Compressing the marked interiors}

\begin{lemma}[Subcritical marked-cube compression]\label{hp:hide-cubes}
Suppose $2<p<3$, and let $U$ and $w_U$ be the marked cubes and
collars of Proposition~\ref{hp:carrier}.  Fix any additional error
$\varepsilon_0>0$.  Let
$k:\Omega\to\Lambda$ be a locally bi-Lipschitz homeomorphism.
Suppose on thinner collars $C_U'(w_U)$, containing a fixed
relative share of the outer half-collar, one has
\begin{equation}\label{hp:k-collar-assumption}
 \sum_U\frac{\ell_U}{w_U}
        \int_{C_U'(w_U)}|Dk|^p\leq b.
\end{equation}
There are disjoint slightly enlarged cubes $\widetilde U\supset U$,
whose boundaries lie in these outer collars, and source
self-homeomorphisms supported in their closures, such that the
resulting locally bi-Lipschitz map $\widetilde k$ agrees with $k$
outside $\bigcup\widetilde U$ and satisfies
\begin{equation}\label{hp:hidden-cost}
 \sum_U\int_{\widetilde U}|D\widetilde k|^p
                \leq C_p b+\varepsilon_0.
\end{equation}
Its value change on each $\widetilde U$ takes place inside
$k(\widetilde U)$.  Thus previously required fine compact value
accuracy is preserved by choosing the marked cubes sufficiently
fine.  In particular, if the costs in
\eqref{hp:k-collar-assumption} are bounded by a fixed multiple of
\eqref{hp:marked-collar-budget} plus errors summable after
multiplication by $\ell_U/w_U$, all marked-interior costs can be
made arbitrarily small.  No convergence to a fixed boundary trace
is assumed.
\end{lemma}

\begin{proof}
Slice the outer collar of each $U$ by concentric cube boundaries.
The coarea formula for the cube radius and
\eqref{hp:k-collar-assumption} give a slightly larger cube
$\widetilde U$ such that
\begin{equation}\label{hp:chosen-cube-trace}
 \sum_U\ell_U\int_{\partial\widetilde U}|Dk|^p\,d\cH^2
                   \leq Cb.
\end{equation}
Choose its radius also to be a Lebesgue point of the surface integral
as a function of cube radius.  Almost every radius has this property,
so it is compatible with the averaging inequality.  The boundaries
are chosen on the outer side to ensure that all of the previously
exempt interior $U$ is included.  The available disjoint collar
enlargements make the cubes $\widetilde U$ disjoint.

Write a fixed such cube as $x_0+a[-1,1]^3$, and use coordinates
$x=x_0+at\theta$, where $0\leq t\leq1$ and
$\theta\in\partial[-1,1]^3$.  The edges between facets have zero
volume and may be omitted in integration.  For $0<\alpha<1/2$
and $0<\rho<1/2$, let $\Phi$ be the radial self-homeomorphism
with radial function
\[
 \varphi(t)=
 \begin{cases}
 (1-\rho)t/\alpha,&0\leq t\leq\alpha,\\[2pt]
 1-\rho+\rho(t-\alpha)/(1-\alpha),&\alpha\leq t\leq1.
 \end{cases}
\]
It fixes the boundary and is bi-Lipschitz for these fixed positive
parameters.  On the inner cube, $k$ has a finite Lipschitz constant
$L$ on the compact cube, so
\[
 \int_{\{t<\alpha\}}|D(k\circ\Phi)|^p
                  \leq C a^3 L^p\alpha^{3-p}.
\]
Choose $\alpha$ so small that this is below an assigned summable
error.  This is exactly where $p<3$ is used.

On the outer region,
$|D\Phi|\leq C(\varphi(t)/t+\varphi'(t))\leq C/t$.
Cube polar integration therefore gives
\[
 \int_{\{\alpha<t<1\}}|D(k\circ\Phi)|^p
 \leq C_pa^3\int_\alpha^1t^{2-p}
       \int_{\partial[-1,1]^3}
          |Dk(x_0+a\varphi(t)\theta)|^p\,d\cH^2(\theta)\,dt.
\]
For fixed $\alpha$, let $\rho\downarrow0$.  The inner surface
integral is sampled from radii in $[1-\rho,1]$.  The chosen outer
radius is an $L^1$ Lebesgue point.  After the linear substitution
from $t$ to $\varphi(t)$, the weight $t^{2-p}$ is bounded by a
constant depending on the fixed $\alpha$; hence Lebesgue
differentiation applies and the limsup of the last expression is
at most
\[
 C_pa\int_{\partial\widetilde U}|Dk|^p\,d\cH^2
                   \int_\alpha^1t^{2-p}\,dt
 \leq C_p\ell_U\int_{\partial\widetilde U}|Dk|^p\,d\cH^2.
\]
The last constant may depend on $p$, since
$\int_0^1t^{2-p}\,dt=(3-p)^{-1}$.  Take $\rho$ small enough to
add only another assigned summable error.  Applying this procedure
on all cubes and using \eqref{hp:chosen-cube-trace} proves
\eqref{hp:hidden-cost}.

Each source modification fixes its boundary.  Local finiteness and
positive fixed radial slopes therefore preserve local
bi-Lipschitzness and the global homeomorphism onto the same target.
All new values on a cube belong to its old image under $k$.  If
$k$ is finely close to $f$, uniform continuity of $f$ on compact
localizations and the original fineness of the cubes make the
resulting oscillation as small as the required compact accuracy.
The proof sampled the actual ambient derivative of $k$ on a newly
chosen radius, and nowhere used convergence of derivatives on a
previously fixed boundary.
\end{proof}
